% arXiv-ready LaTeX: use standard packages only, no exotic dependencies
\documentclass[12pt,a4paper]{article}

% --- Standard packages (all available on arXiv) ---
\usepackage{amsmath,amssymb,amsthm}
\usepackage[T1]{fontenc}
\usepackage[utf8]{inputenc}


% ISO date: YYYY-MM-DD
\newcommand{\isodatetoday}{%
  \the\year-%
  \ifnum\month<10 0\fi\the\month-%
  \ifnum\day<10 0\fi\the\day}
\newcommand{\paperversion}{v4.15}
\usepackage{geometry}

\usepackage{setspace}
\usepackage{booktabs}
\usepackage{array}
\usepackage{longtable}
\usepackage{enumitem}
\usepackage{xcolor}
\usepackage[colorlinks=true,
            linkcolor=black,
            citecolor=black,
            urlcolor=blue,
            pdftitle={On the Riemann Hypothesis: A Kernel-Theoretic Approach via Tensor Contraction and Self-Adjoint Extension},
            pdfauthor={Marc Brendecke},
            pdfsubject={Mathematics},
            pdfkeywords={Riemann Hypothesis, dynamical zeta function, self-adjoint extension, tensor contraction, Hilbert-Polya, Berry-Keating}]{hyperref}

\geometry{a4paper, top=3cm, bottom=3cm, left=3.5cm, right=3.5cm}
\setstretch{1.15}

% --- Theorem environments ---
\newtheorem{theorem}{Theorem}[section]
\newtheorem{lemma}[theorem]{Lemma}
\newtheorem{proposition}[theorem]{Proposition}
\newtheorem{corollary}[theorem]{Corollary}
\newtheorem{definition}[theorem]{Definition}
\newtheorem{axiom}[theorem]{Axiom}
\newtheorem{assumption}[theorem]{Assumption}
\theoremstyle{remark}
\newtheorem*{remark}{Remark}

% --- Custom commands ---
\newcommand{\calH}{\mathcal{H}}
\newcommand{\calZ}{\mathcal{Z}}
\newcommand{\calL}{\mathcal{L}}
\newcommand{\calP}{\mathcal{P}}
\newcommand{\calC}{\mathcal{C}}
\newcommand{\calJ}{\mathcal{J}}
\newcommand{\half}{\tfrac{1}{2}}
\newcommand{\Adj}{\mathrm{Adj}}

% ====================================================================
\begin{document}
% ====================================================================

\title{\textbf{On the Riemann Hypothesis:}\\[4pt]
A Kernel-Theoretic Approach via\\
Tensor Contraction and Self-Adjoint Extension}

\author{Marc Brendecke\\[4pt]
\small Independent Researcher\\
\small Bochum, Germany\\[2pt]
\small \texttt{[contact on request]}}

\date{\isodatetoday\\[2pt]
\small\textnormal{Version \paperversion\ --- Preprint, not peer reviewed}\\[4pt]
\small\textnormal{DOI: \href{https://doi.org/10.5281/zenodo.19157158}{10.5281/zenodo.19157158}}\\[4pt]
\small\textnormal{\copyright\ \the\year\ Marc Brendecke. All rights reserved.}}

\maketitle
\thispagestyle{empty}

\begin{abstract}
\noindent
We develop a conditional proof of the Riemann Hypothesis (RH) from
a minimal combinatorial foundation.
Starting from a kernel $(S,\Phi)$---a set equipped with a
self-map---eleven lemmas establish the orbit structure and yield,
via a corollary (K1) and a Laplace identity (L11), the dynamical zeta
function $\mathcal{Z}(s)$.
Identifying the kernel with the arithmetic kernel
$K_{\mathrm{arith}}=(\mathbb{P},\mathrm{id})$ with weight
$\lambda(\{p\})=\log p$ gives $\mathcal{Z}(s)=\zeta(s)$.

The functional equation $\xi(s)=\xi(1-s)$ is interpreted as the
vanishing of an adjunction tensor $\mathrm{Adj}_J(\mathcal{L}_0)=0$,
where $J:s\mapsto1-\bar{s}$ is an anti-unitary involution on the
Hardy space $\mathcal{H}^+$.
The Berry--Keating operator $\mathcal{L}_0=-i(xd/dx+\frac{1}{2})$
is shown to be essentially self-adjoint on the $J$-invariant subspace
$\mathcal{H}^+_+$ via a Levi-Civita tensor argument
($\varepsilon^{\mu\nu}$-annihilation of deficiency elements).

The Heisenberg compensator $C(f)=\frac{1}{2}(f+Jf)$ is introduced as
the map projecting any deficiency element to zero; its tensor form
$\mathcal{C}^{\mu}{}_{\nu}$ is shown to be idempotent and is identified
with the cascade operator of the Sphaera-Cascade framework.

A discreteness lemma for the spectrum of $\mathcal{L}_0$ on
$\mathcal{H}^+_+$, grounded in the prime lattice acting as a compact
quotient, closes the argument conditionally on an identification with
Connes' adèle class space.

RH then follows as the vanishing of the contravariant spectral
component: $\gamma_\rho^-=0$ for all non-trivial zeros $\rho$,
i.e.\ $\operatorname{Re}(\rho)=\frac{1}{2}$.

\medskip
\noindent\textit{MSC 2020:} 11M26 (primary); 47B25, 58J50, 11M06 (secondary).

\noindent\textit{Keywords:} Riemann Hypothesis, dynamical zeta function,
self-adjoint extension, tensor contraction, Hilbert--P\'olya conjecture,
Berry--Keating operator, Heisenberg compensator.
\end{abstract}

\newpage
\tableofcontents
\newpage

% ====================================================================
\section*{Nomenclature}
\addcontentsline{toc}{section}{Nomenclature}
% ====================================================================

\begin{longtable}{p{3.2cm}p{11cm}}
\toprule
\textbf{Symbol} & \textbf{Meaning}\\
\midrule
\endfirsthead
\multicolumn{2}{l}{\textit{(Nomenclature continued)}}\\[2pt]
\toprule
\textbf{Symbol} & \textbf{Meaning}\\
\midrule
\endhead
\midrule
\multicolumn{2}{r}{\textit{(continued on next page)}}\\
\endfoot
\bottomrule
\endlastfoot
%
\multicolumn{2}{l}{\textit{Sets, maps and orbit structure}}\\[2pt]
$S$, $\Phi$ & Carrier set and self-map of the dynamical system\\
$\mathbb{P}$ & Set of prime numbers\\
$\Gamma_\mathbb{P}$ & Prime lattice $\langle p:p\in\mathbb{P}\rangle\subset\mathbb{Q}^+$\\
$\mathbb{A}$, $\mathbb{A}^1$ & Adèle ring of $\mathbb{Q}$; norm-$1$ idèles $\{x\in\mathbb{A}^\times:|x|_\mathbb{A}=1\}$\\
$O(x)$, $R(x)$ & Forward orbit; return set $\{n\ge1:\Phi^n(x)=x\}$\\
$\ell(x)$, $\ell(P)$ & Minimal period of $x$; combinatorial orbit length\\
$P(x)$, $\mathcal{P}$ & Primitive orbit through $x$; collection of all primitive orbits\\
$\lambda(P)$, $h$ & Geometric weight (Axiom~W); growth exponent (Axiom~G)\\
$K_{\mathrm{arith}}$ & Arithmetic kernel $(\mathbb{P},\mathrm{id})$ with $\lambda(\{p\})=\log p$\\
$K_\pi$ & Automorphic kernel for self-dual $\pi$ (same $(\mathbb{P},\mathrm{id},\log p)$)\\[4pt]
%
\multicolumn{2}{l}{\textit{Analytic objects}}\\[2pt]
$\Theta(\tau)\to\Theta^\tau$ & Orbit distribution (L9); promoted to contravariant vector at L11\\
$\mathcal{Z}(s)$ & Dynamical zeta function $\prod_P(1-e^{-s\lambda(P)})^{-1}$\\
$\zeta(s)$, $\xi(s)$ & Riemann zeta; completed zeta $\tfrac{1}{2}s(s-1)\pi^{-s/2}\Gamma(\tfrac{s}{2})\zeta(s)$\\
$\rho=\tfrac{1}{2}+i\gamma_\rho$ & Non-trivial zero; spectral parameter\\
$\Sigma(\tau)$, $R(\tau)$ & Spectral distribution; trivial-zero term $-e^{-\tau/2}/(e^\tau-1)$\\[4pt]
%
\multicolumn{2}{l}{\textit{Tensor notation}}\\[2pt]
$\mu,\nu\in\{+,-\}$ & Covariant / contravariant sector indices\\
$z^+=\tfrac{1}{2}$, $z^-=z-\tfrac{1}{2}$ & Sector components; RH $\Leftrightarrow\operatorname{Re}(\gamma_\rho^-)=0$\\
$\varepsilon^{\mu\nu}$ & 2D Levi-Civita tensor ($\varepsilon^{+-}=1$, $\varepsilon^{-+}=-1$)\\
$\mathcal{L}^s{}_\tau=e^{-s\tau}$ & Laplace kernel (mixed tensor, L11)\\
$J\to\mathcal{J}^{\mu}{}_\nu$ & Reflection $s\mapsto1-\bar{s}$; sector matrix $\mathrm{diag}(1,-1)$\\
$U(\tau)\to U^\mu{}_\nu$ & Propagator $e^{-i\tau\mathcal{L}_0}$; trace $U^\mu{}_\mu=\Sigma(\tau)$\\
$\mathrm{Adj}_J(\mathcal{O})$ & Adjunction tensor $(J\mathcal{O}J^{-1}-\mathcal{O})$\\[4pt]
%
\multicolumn{2}{l}{\textit{Operators and spaces}}\\[2pt]
$\mathcal{H}^+$, $\mathcal{H}^+_+$ & Hardy space $H^2(\operatorname{Re}(s)>\tfrac{1}{2})$; $J$-invariant subspace\\
$\mathcal{H}_0$ & Weighted space $L^2(\mathbb{R}^+,e^{-x}dx)$\\
$J$ & Anti-unitary involution $s\mapsto1-\bar{s}$\\
$\mathcal{L}_0$ & Berry--Keating operator $-i(xd/dx+\tfrac{1}{2})$ on $\mathcal{H}_0$\\
$(n_+,n_-)$ & Deficiency indices; $(0,0)$ on $\mathcal{H}^+_+$ (Theorem~\ref{thm:T45})\\
$C$, $\mathcal{C}^{\mu}{}_\nu$ & Heisenberg compensator map and tensor (\S\ref{sec:compensator})\\[4pt]
%
\multicolumn{2}{l}{\textit{Sphaera-Cascade framework}}\\[2pt]
$\mathbb{S}$ & Sphaera: maximal total set; realized as $\mathcal{H}^+_+$\\
$S_n$, $S_0$ & Strata of the cascade; terminal ground state $\{0\}$\\
$\hat{K}$ & Cascade operator; realized as $\mathcal{C}^\mu{}_\nu\circ\Pi_J$\\
$\Pi_J$ & $J$-symmetrization $f\mapsto\tfrac{1}{2}(f+Jf)$; projects $\mathcal{H}^+\to\mathcal{H}^+_+$\\
K1 & Kernel corollary: $\mathcal{Z}$ encodes orbit structure\\
L11 & Index contraction: $\hat\Theta^s=\log\mathcal{Z}(s)$\\
T1+T3 & $\mathrm{Adj}_J(\mathcal{L}_0)=0$ (formal self-adjointness)\\
T4+T5 & $\varepsilon^{\mu\nu}$-annihilation; deficiency indices $(0,0)$\\
T6 & Trace formula $U^\mu{}_\mu=\Sigma(\tau)$\\[4pt]
%
\multicolumn{2}{l}{\textit{Dixon-algebraic hierarchy (\S\ref{sec:dixon}--\ref{sec:septonian})}}\\[2pt]
$\mathcal{B}=\mathbb{C}\otimes\mathbb{H}$ & Biquaternion algebra ($\dim_\mathbb{R}=8$)\\
$e_1,e_2,e_3$ & Quaternion imaginary units; $e_ae_b=\varepsilon_{abc}e_c$\\
$J_{\mathcal{B}}$, $R_{e_1}$ & Lift of $J$ to $\mathcal{B}$; $e_1$-reflection in $\mathbb{H}$\\
$\iota(\rho)$ & Embedding $\rho\mapsto\tfrac{1}{2}+(\sigma-\tfrac{1}{2})e_1+\gamma i\in\mathcal{B}$\\
$\mathcal{D}=\mathbb{C}\otimes\mathbb{O}$ & Octo-Dixon algebra ($\dim_\mathbb{R}=16$)\\
$\mathbb{O}$, $\ell$ & Octonions; first Cayley-Dickson doubling unit ($\mathbb{O}=\mathbb{H}\oplus\mathbb{H}\ell$)\\
$\mathcal{T}_1,\mathcal{T}_2,\mathcal{T}_3$ & Three Trioles in $\mathbb{O}$; incidence governed by Fano plane $\mathrm{PG}(2,2)$\\
$\Omega$ & General Octo-Dixon element $i(t\cdot1_\mathbb{O}+\sum_{a=1}^7 q_ae_a)$\\
$J_{\mathcal{D}}$ & $J$-involution on $\mathcal{D}$; negates $\{e_1,e_1\ell\}$\\
$\mathbb{S}_{\mathrm{ed}}$, $\ell'$ & Sedenions (CD double of $\mathbb{O}$, $\dim=16$, 15 imaginaries); 2nd CD unit\\
$f_1,\ldots,f_{15}$ & Sedenion imaginary units ($f_a:=e_a$ for $a\le7$; $f_{8+k}:=e_k\ell'$)\\
$\mathcal{O}_1,\mathcal{O}_2,\mathcal{O}_3$ & Three Octo-Trioles in $\mathbb{S}_{\mathrm{ed}}$; governed by $\mathrm{PG}(3,2)$\\
$\mathcal{S}_{\mathrm{ep}}=\mathbb{C}\otimes\mathbb{S}_{\mathrm{ed}}$ & Septonian algebra ($\dim_\mathbb{R}=32$)\\
$\Psi$ & General Septonian element $i(t\cdot1_{\mathbb{S}_{\mathrm{ed}}}+\sum_{a=1}^{15}r_af_a)$\\
$J_{\mathcal{S}}$ & $J$-involution on $\mathcal{S}_{\mathrm{ep}}$; negates $\{f_1,f_5,f_9,f_{13}\}=\{e_1,e_1\ell,e_1\ell',e_1\ell\ell'\}$\\
$\mathrm{PG}(n,2)$ & Projective space over $\mathbb{F}_2$; governs Triole incidence at level $n$\\[4pt]
%
\multicolumn{2}{l}{\textit{Field equation and probe structure (\S\ref{sec:field_probe})}}\\[2pt]
$P_+$, $P_-$ & Sector projectors $\tfrac{1}{2}(1\pm J)$; onto $\mathcal{H}^+_+$ resp.\ $\mathcal{H}^+_-$\\
$\Pi_L$, $\Pi_R$ & Left/right probe operators; orthogonal ($\Pi_L\Pi_R=0$)\\
$\mathcal{G}_L$, $\mathcal{G}_R$ & Guard subspaces $\mathrm{ran}(\Pi_L)$, $\mathrm{ran}(\Pi_R)$\\
$\mathcal{H}_{\mathrm{search}}$ & Search space $(\mathcal{G}_L\oplus\mathcal{G}_R)^\perp$\\
$L_\perp$ & Compression $P_\perp\mathcal{L}_0 P_\perp$ of $\mathcal{L}_0$ to $\mathcal{H}_{\mathrm{search}}$\\
\end{longtable}

\newpage

% ====================================================================
% !! GESPERRT: Abschnitt 1 (Kombinatorische Orbittheorie) !!
% D1, L1-L6, Axiom W, Axiom G, L7-L11, K1
% Diese Lemmas sind mengenunabhaengig: sie funktionieren fuer jede
% Menge S mit beliebiger Selbstabbildung Phi.
% Keine Aenderungen ohne ausdrueckliche Freigabe.
% Geprueft und gesperrt: v4, 2026-03-20
% ====================================================================
\section{Combinatorial Orbit Theory}
\label{sec:combinatorial}
% ====================================================================

\begin{definition}[Dynamical system]\label{def:D1}
A \emph{dynamical system} is a pair $(S,\Phi)$ where $S$ is an
arbitrary set and $\Phi:S\to S$ an arbitrary map.
No topology, metric, or measure on $S$ is assumed.
The collection of primitive orbits $\mathcal{P}(S,\Phi)$ is defined
through Lemmas \ref{lem:L1}--\ref{lem:L6}.
\end{definition}

\begin{lemma}\label{lem:L1}
$x\in S$ is \emph{periodic} if and only if the return set
$R(x):=\{n\ge1:\Phi^n(x)=x\}$ is nonempty.
\end{lemma}
\begin{proof}By definition.\end{proof}

\begin{remark}[Notation]
For periodic $x$, $R(x)\subseteq\mathbb{N}_{\ge1}$ is nonempty,
so the following minimization is well-defined.
\end{remark}

\begin{lemma}\label{lem:L3}
Every periodic $x$ has a unique minimal period
$\ell(x):=\min R(x)\in\mathbb{N}_{\ge1}$.
\end{lemma}
\begin{proof}The well-ordering principle applied to the nonempty set $R(x)\subseteq\mathbb{N}$.\end{proof}

\begin{lemma}\label{lem:L4}
$\ell(x)=k$ if and only if $\Phi^k(x)=x$ and $\Phi^m(x)\neq x$ for all $1\le m<k$.
The orbit $\{x,\Phi(x),\ldots,\Phi^{k-1}(x)\}$ consists of exactly $k$ distinct points.
\end{lemma}
\begin{proof}
The equivalence follows from the definition of $\min$.
Distinctness: if $\Phi^i(x)=\Phi^j(x)$ for $0\le i<j\le k-1$,
then $\Phi$ is invertible on the finite orbit, giving
$\Phi^{j-i}(x)=x$ with $j-i<k$, contradicting minimality.
\end{proof}

\begin{definition}[Primitive orbit]
The \emph{primitive orbit} through periodic $x$ with $\ell(x)=k$ is
$P(x):=\{x,\Phi(x),\ldots,\Phi^{k-1}(x)\}$, with
\emph{combinatorial length} $\ell(P):=k$.
Write $\mathcal{P}=\mathcal{P}(S,\Phi)$ for the collection of all primitive orbits.
\end{definition}

\begin{lemma}\label{lem:L5}
Every periodic point belongs to exactly one primitive orbit $P(x)$,
and $P(\Phi^j(x))=P(x)$ for all $j\ge0$.
\end{lemma}
\begin{proof}
\emph{Generation.} For $m\ge0$ write $m=qk+r$, $0\le r<k$; then
$\Phi^m(x)=\Phi^r(x)\in P(x)$, so $O(x)=P(x)$.
\emph{Uniqueness.} If $Q\ni x$ is primitive then
$Q\subseteq O(x)=P(x)$ and $P(x)\subseteq Q$, hence $Q=P(x)$.
\emph{Consistency.} $\Phi^k(\Phi^j(x))=\Phi^j(x)$ gives
$\ell(\Phi^j(x))\le k$; if $\ell(\Phi^j(x))=k'<k$ then
$\Phi^{k'}(x)=x$, contradicting $k=\ell(x)$.
\end{proof}

\begin{lemma}\label{lem:L6}
$\mathcal{P}$ is pairwise disjoint and covers all periodic points.
\end{lemma}
\begin{proof}
If $y\in P\cap Q$ then $P=P(y)=Q$ by Lemma~\ref{lem:L5}.
Every periodic $x$ lies in $P(x)\in\mathcal{P}$.
\end{proof}

\begin{axiom}[W -- Weight function]\label{ax:W}
A map $\lambda:\mathcal{P}\to\mathbb{R}_{>0}$ assigns a positive real
\emph{geometric length} to each primitive orbit.
No relation between $\lambda(P)$ and $\ell(P)$ is assumed.
\end{axiom}

\begin{axiom}[G -- Growth condition]\label{ax:G}
There exists $h\ge0$ such that
$\#\{P\in\mathcal{P}:\lambda(P)\le T\}=O(e^{hT})$ as $T\to\infty$.
\end{axiom}

\begin{remark}
Axiom~W is stated here, after $\mathcal{P}$ is defined (Lemma~\ref{lem:L6}),
because $\lambda$ must be defined on primitive orbits.
Axiom~G is stated immediately after Axiom~W because it references both
$\mathcal{P}$ and $\lambda$; placing it before either would be a
forward reference.
Axioms W and G are the only analytic hypotheses;
Lemmas \ref{lem:L1}--\ref{lem:L6} are purely combinatorial.
\end{remark}

\begin{lemma}\label{lem:L7}
The $m$-fold cover of $P\in\mathcal{P}$ with $\ell(P)=k$ is a closed
path of length $mk$ with exactly $m$ cyclic symmetries and no others.
\end{lemma}
\begin{proof}
Rotation by $r$ positions is a symmetry iff $k\mid r$.
The multiples of $k$ in $\{0,\ldots,mk-1\}$ number exactly $m$.
\end{proof}

\begin{lemma}\label{lem:L8}
If $x$ is periodic and $n\in R(x)$, then $\ell(x)\mid n$.
\end{lemma}
\begin{proof}
Write $n=q\ell(x)+r$, $0\le r<\ell(x)$. Then $\Phi^r(x)=x$,
so $r=0$ by minimality of $\ell(x)$.
\end{proof}

\begin{lemma}[Orbit measure]\label{lem:L9}
\[
\Theta(\tau)
:= \sum_{P\in\mathcal{P}}\sum_{m\ge1}
\frac{1}{m}\,\delta\!\left(\tau-m\lambda(P)\right)
\]
is the weighted distribution of closed orbit lengths.
\end{lemma}
\begin{proof}
By Lemma~\ref{lem:L8}, every closed orbit of geometric length $\tau$
is an $m$-fold cover of a unique $P$ with $\lambda(P)=\tau/m$.
By Lemma~\ref{lem:L7} it carries exactly $m$ cyclic symmetries,
contributing weight $1/m$.
\end{proof}

\begin{remark}
The promotion $\Theta(\tau)\to\Theta^\tau$ (distribution to
contravariant vector) occurs at Lemma~\ref{lem:L11}, where $\tau$ becomes
a contracted index.
\end{remark}

\begin{lemma}[Dynamical zeta function]\label{lem:L10}
For $\operatorname{Re}(s)>h$ the product
\[
\mathcal{Z}(s):=\prod_{P\in\mathcal{P}}\!\left(1-e^{-s\lambda(P)}\right)^{-1}
\]
converges absolutely and defines a holomorphic nonvanishing function.
\end{lemma}
\begin{proof}
Absolute convergence of $\sum_P e^{-\sigma\lambda(P)}$ for $\sigma>h$
follows from Axiom~G by grouping orbits into dyadic length blocks.
\end{proof}

\begin{corollary}[K1 --- Kernel corollary]\label{cor:K1}
The dynamical system $(S,\Phi,\lambda)$ canonically produces
$\mathcal{Z}(s)$, which encodes the primitive orbit structure completely.
The internal structure of $S$ is irrelevant: only $(\mathcal{P},\lambda)$ determine $\mathcal{Z}$.
\end{corollary}

\begin{lemma}[L11 --- Index contraction]\label{lem:L11}
\[
\hat\Theta^s
\;:=\; \mathcal{L}^s{}_\tau\cdot\Theta^\tau
\;=\; \int_0^\infty e^{-s\tau}\,\Theta(\tau)\,d\tau
\;=\; \log\mathcal{Z}(s),
\qquad \operatorname{Re}(s)>h.
\]
\end{lemma}
\begin{proof}
From Corollary~\ref{cor:K1}: $\log\mathcal{Z}(s)=-\sum_P\log(1-e^{-s\lambda(P)})$.
From Lemma~\ref{lem:L9}: the left side equals
$\sum_P\sum_{m\ge1}\frac{1}{m}e^{-sm\lambda(P)}$.
Applying $\sum_{m\ge1}\frac{x^m}{m}=-\log(1-x)$ and using absolute
convergence from Lemma~\ref{lem:L10} to justify term-by-term integration
gives the result.
\end{proof}

% ====================================================================
% !! ENDE GESPERRTER ABSCHNITT !!
% ====================================================================

% ====================================================================
% !! GESPERRT: Abschnitt 2 (Arithmetische Identifikation) !!
% Prop. arith: Axiom G mit h=1, Z(s)=zeta(s), von-Mangoldt-Formel.
% Klassische Arithmetik, korrekt geprueft.
% Keine Aenderungen ohne ausdrueckliche Freigabe.
% Geprueft und gesperrt: v4, 2026-03-20
% ====================================================================
\section{Arithmetic Identification}
\label{sec:arithmetic}
% ====================================================================

\begin{definition}[Arithmetic kernel]
$K_{\mathrm{arith}}:=(\mathbb{P},\mathrm{id}_\mathbb{P})$
with weight $\lambda(\{p\}):=\log p$.
Every prime is a fixed point; primitive orbits are singletons.
\end{definition}

\begin{proposition}\label{prop:arith}
\begin{enumerate}[label=(\roman*),leftmargin=*]
\item Axiom~G holds with $h=1$: $\pi(e^T)\le e^T$.
\item $e^{-s\log p}=p^{-s}$, so $\mathcal{Z}(s)=\prod_p(1-p^{-s})^{-1}=\zeta(s)$
for $\operatorname{Re}(s)>1$.
\item Lemma~\ref{lem:L11} reproduces the von~Mangoldt formula
$\log\zeta(s)=\int_0^\infty e^{-s\tau}\Theta(\tau)\,d\tau$.
\end{enumerate}
The arithmetic enters here and is not re-invoked until
Lemma~\ref{lem:discrete}.
\end{proposition}

% ====================================================================
% !! ENDE GESPERRTER ABSCHNITT (Arithmetische Identifikation) !!
% ====================================================================

% ====================================================================
% !! GESPERRT: Abschnitt 3 (Meromorphe Fortsetzung und Funktionalgleichung) !!
% Thm. FE: xi(s)=xi(1-s), klassisch via Jacobi-Theta + Poisson.
% Thm. EF: explizite Formel Guinand-Weil, korrekte Skizze.
% Keine Aenderungen ohne ausdrueckliche Freigabe.
% Geprueft und gesperrt: v4, 2026-03-20
% ====================================================================
\section{Meromorphic Continuation and Functional Equation}
\label{sec:FE}
% ====================================================================

\begin{definition}[Completed zeta function]
$\xi(s):=\tfrac{1}{2}s(s-1)\pi^{-s/2}\Gamma\!\left(\tfrac{s}{2}\right)\zeta(s)$.
\end{definition}

\begin{theorem}[Functional equation]\label{thm:FE}
$\xi$ extends to an entire function satisfying $\xi(s)=\xi(1-s)$.
$\zeta(s)$ is meromorphic on $\mathbb{C}$ with a simple pole at $s=1$;
all non-trivial zeros lie in the critical strip $0<\operatorname{Re}(s)<1$.
\end{theorem}
\begin{proof}[Proof (sketch)]
The Jacobi theta function $\vartheta(t)=\sum_n e^{-\pi n^2 t}$
satisfies $\vartheta(t)=t^{-1/2}\vartheta(1/t)$ by Poisson summation.
The Mellin representation
$\pi^{-s/2}\Gamma(\tfrac{s}{2})\zeta(s)
=\int_0^\infty t^{s/2-1}\tfrac{\vartheta(t)-1}{2}\,dt$,
split at $t=1$ and symmetrized, yields $\xi(s)=\xi(1-s)$.
\end{proof}

\begin{theorem}[Explicit formula --- Guinand--Weil]\label{thm:EF}
For $\tau>0$, in the distributional sense:
\[
\Theta(\tau) = 1 - \Sigma(\tau) + R(\tau),
\qquad
\Sigma(\tau):=\sum_\rho e^{-i\tau\gamma_\rho},
\qquad
R(\tau):=\frac{-e^{-\tau/2}}{e^\tau-1}.
\]
\end{theorem}
\begin{proof}[Proof (sketch)]
Apply the inverse Mellin transform to $\log\zeta(s)=\hat\Theta^s$
(Lemma~\ref{lem:L11}). Residues at $s=1$ contribute $+1$;
non-trivial zeros $\rho$ contribute $-\Sigma$;
trivial zeros $-2n$ via the Hadamard product of $\xi$ contribute $R$.
\end{proof}

% ====================================================================
% !! ENDE GESPERRTER ABSCHNITT (Funktionalgleichung) !!
% ====================================================================

% ====================================================================
% !! GESPERRT: Abschnitt 4 (J-Tensor und Sektorzerlegung) !!
% Def. J: anti-unitaere Involution, H^+_+ = ker(J-id), RH-Aequivalenz.
% Korrekt und praezise.
% Keine Aenderungen ohne ausdrueckliche Freigabe.
% Geprueft und gesperrt: v4, 2026-03-20
% ====================================================================
\section{The $J$-Tensor and Sector Decomposition}
\label{sec:J}
% ====================================================================

\begin{definition}[$J$-tensor and sector decomposition]\label{def:J}
The reflection $J:s\mapsto1-\bar{s}$ is an anti-unitary involution
on the Hardy space
$\mathcal{H}^+:=H^2(\{\operatorname{Re}(s)>\tfrac{1}{2}\})$.
For any $z\in\mathbb{C}$ define the \emph{covariant} and
\emph{contravariant} components:
\[
z^+ := \tfrac{1}{2}, \qquad z^- := z-\tfrac{1}{2}.
\]
Then $z=z^++z^-$ and $\operatorname{Re}(z)=\tfrac{1}{2}$ iff
$\operatorname{Re}(z^-)=0$.
The $J$-invariant subspace is
$\mathcal{H}^+_+:=\{f\in\mathcal{H}^+:Jf=f\}$.
\end{definition}

\begin{remark}[Promotion of $J$ to $\mathcal{J}^{\mu}{}_\nu$]
In the sector basis $\mu,\nu\in\{+,-\}$, with
$z=z^+e_++z^-e_-$, the map $J$ acts as $z^+\mapsto z^+$
and $z^-\mapsto-\overline{z^-}$ (for real $z^-$: sign flip).
Its matrix representation is
\[
\mathcal{J}^{\mu}{}_\nu
:= \begin{pmatrix}1&0\\0&-1\end{pmatrix}.
\]
The subspace $\mathcal{H}^+_+$ is $\ker(\mathcal{J}-\mathbf{1})$,
i.e.\ functions with vanishing contravariant component.
\textbf{RH} is equivalent to: all $\gamma_\rho$ lie in $\ker(\mathcal{J}-\mathbf{1})$.
\end{remark}

% ====================================================================
% !! ENDE GESPERRTER ABSCHNITT (J-Tensor) !!
% ====================================================================

% ====================================================================
% !! GESPERRT: Abschnitt 4a (Dixon-Algebra) !!
% Biquaternionische Realisierung der Sektorzerlegung.
% Prop. JB: J_B = e1-Reflexion, liftet J exakt auf iota. BEWIESEN.
% Kor. RH_quaternion: RH <=> e1-Komponente=0. BEWIESEN.
% Prop. shadow: 2x2-Kalkuel als Projektion auf {1,e1}-Ebene. BEWIESEN.
% Rem. triole: S^3, Hopf-Faserung, Fixpunkt = kritische Linie. KORREKT.
% Keine Aenderungen ohne ausdrueckliche Freigabe.
% Geprueft und gesperrt: v4_12, 2026-03-22
% ====================================================================
\section{Dixon-Algebraic Realization of the Sector Structure}
\label{sec:dixon}
% ====================================================================

The $2\times2$ tensor calculus of \S\ref{sec:J} operates with a single
contravariant direction $e_-$. 
We show here that this structure is the projection of a richer
biquaternion algebra, which reveals three orthogonal ``contravariant
channels'' corresponding to the Triole of the Sphaera framework.

\subsection*{The Biquaternion Algebra}

\begin{definition}[Biquaternion algebra $\mathcal{B}$]\label{def:biquaternion}
The \emph{biquaternion algebra} is $\mathcal{B}:=\mathbb{C}\otimes_{\mathbb{R}}\mathbb{H}$,
with basis $\{1,\,e_1,\,e_2,\,e_3,\,i,\,ie_1,\,ie_2,\,ie_3\}$
over $\mathbb{R}$ ($\dim_{\mathbb{R}}\mathcal{B}=8$).
Here $i\in\mathbb{C}$ satisfies $i^2=-1$ and commutes with all
$e_a\in\mathbb{H}$; the quaternion units obey
$e_a^2=-1$ and $e_ae_b=\varepsilon_{abc}e_c$ for $a\neq b$.

The \emph{biquaternion conjugate} of
$q=\alpha+\beta_1e_1+\beta_2e_2+\beta_3e_3
   +(\gamma+\delta_1e_1+\delta_2e_2+\delta_3e_3)\,i$
is
$\bar{q}:=\alpha-\beta_1e_1-\beta_2e_2-\beta_3e_3
   +(\gamma-\delta_1e_1-\delta_2e_2-\delta_3e_3)\,i$.
\end{definition}

\begin{definition}[Embedding of $\mathbb{C}$ into $\mathcal{B}$
and Triole decomposition]\label{def:triole_embed}
For a non-trivial zero $\rho=\sigma+i\gamma$ of $\zeta$, define
\[
  \iota(\rho)
  \;:=\;
  \underbrace{\tfrac{1}{2}\cdot 1}_{\text{covariant sector }e_+}
  \;+\;
  \underbrace{(\sigma-\tfrac{1}{2})\cdot e_1}_{\text{off-critical channel}}
  \;+\;
  \gamma\cdot i
  \;\in\;\mathcal{B}.
\]
The three mutually orthogonal \emph{Triole} directions in $\mathbb{H}$
are $\mathbf{T}:=(e_1,e_2,e_3)$.
Each pair satisfies $e_ae_b+e_be_a=-2\delta_{ab}$ (mutual
anti-commutation, hence orthogonality in the natural inner product
on $\mathrm{Im}(\mathbb{H})\cong\mathbb{R}^3$).
The full \emph{Triole basis} of $\mathcal{B}$ is
$(1,\;e_1i,\;e_2i,\;e_3i) = (1,\;ie_1,\;ie_2,\;ie_3)$.
\end{definition}

\subsection*{The $J$-Involution as a Biquaternion Reflection}

\begin{proposition}[Biquaternion realization of $J$]
\label{prop:JB}
Define the linear map $J_{\mathcal{B}}:\mathcal{B}\to\mathcal{B}$ by
\[
  J_{\mathcal{B}}(q + p\,i)
  \;:=\;
  \bigl(R_{e_1}(q)\bigr) + p\,i,
\]
where $R_{e_1}:\mathbb{H}\to\mathbb{H}$ is the \emph{reflection}
negating the $e_1$-component and fixing $1,e_2,e_3$:
\[
  R_{e_1}(a + b_1e_1 + b_2e_2 + b_3e_3)
  \;:=\; a - b_1e_1 + b_2e_2 + b_3e_3.
\]
Then $J_{\mathcal{B}}\circ\iota = \iota\circ J$, i.e.\ $J_{\mathcal{B}}$
lifts the anti-unitary involution $J:s\mapsto1-\bar{s}$ to $\mathcal{B}$.
\end{proposition}
\begin{proof}
Let $\rho=\sigma+i\gamma$.
Then $J(\rho)=1-\overline{\rho}=(1-\sigma)+i\gamma$.
Under $\iota$:
\[
  \iota(J(\rho))
  = \tfrac{1}{2}\cdot 1 + (\tfrac{1}{2}-\sigma)\,e_1 + \gamma\,i.
\]
On the other hand, 
\[
  J_{\mathcal{B}}(\iota(\rho))
  = J_{\mathcal{B}}\!\bigl(\tfrac{1}{2} + (\sigma-\tfrac{1}{2})e_1 + \gamma i\bigr)
  = \tfrac{1}{2}\cdot 1 - (\sigma-\tfrac{1}{2})e_1 + \gamma\,i
  = \iota(J(\rho)).
\qedhere
\]
\end{proof}

\begin{corollary}[Fixed-point space and RH in $\mathcal{B}$]
\label{cor:RH_quaternion}
The fixed-point space of $J_{\mathcal{B}}$ on $\iota(\mathbb{C})$ is
\[
  \mathrm{Fix}(J_{\mathcal{B}})
  \;=\;
  \bigl\{\,q\in\iota(\mathbb{C})\;:\;e_1\text{-coefficient of }q=0\,\bigr\}
  \;=\;
  \bigl\{\,\tfrac{1}{2} + \gamma\,i : \gamma\in\mathbb{R}\,\bigr\}.
\]
Therefore the Riemann Hypothesis is equivalent to:
\[
  \mathrm{RH}
  \;\Longleftrightarrow\;
  \iota(\rho)\in\mathrm{Fix}(J_{\mathcal{B}})
  \quad\text{for all non-trivial zeros }\rho,
\]
i.e.\ the $e_1$-component of every $\iota(\rho)$ vanishes.
\end{corollary}

\subsection*{The 2$\times$2 Calculus as a Projection}

\begin{proposition}[2$\times$2 sector calculus as $\mathcal{B}$-shadow]
\label{prop:shadow}
Let $\pi_{01}:\mathcal{B}\to\mathrm{span}_{\mathbb{R}}\{1,e_1\}$
be the projection discarding $e_2$-, $e_3$-, and all $i$-components.
Under the identification
\[
  1\;\longleftrightarrow\; e_+,
  \qquad
  e_1\;\longleftrightarrow\; e_-,
\]
the following holds:
\begin{enumerate}[label=(\roman*),leftmargin=*]
\item $\pi_{01}\circ J_{\mathcal{B}}\circ\pi_{01}^{-1}
      = \mathcal{J}^{\mu}{}_\nu = \mathrm{diag}(1,-1)$.
\item $\pi_{01}\circ R_{e_1}\circ\pi_{01}^{-1}
      = \mathcal{C}^{\mu}{}_\nu = \bigl(\begin{smallmatrix}1&1\\0&0\end{smallmatrix}\bigr)$.
\end{enumerate}
Hence the $2\times2$ tensor calculus of \S\ref{sec:J}--\ref{sec:compensator}
is the restriction of the biquaternion structure to the $\{1,e_1\}$-plane.
\end{proposition}
\begin{proof}
(i) follows directly from Proposition~\ref{prop:JB}: $R_{e_1}(1)=1$,
$R_{e_1}(e_1)=-e_1$, which under $e_+\leftrightarrow1$,
$e_-\leftrightarrow e_1$ gives $\mathcal{J}=\mathrm{diag}(1,-1)$.\\
(ii) The projection onto $\mathrm{Fix}(R_{e_1})=\mathrm{span}\{1,e_2,e_3\}$
restricted to $\mathrm{span}\{1,e_1\}$ maps $a+be_1\mapsto a$.
In the sector basis: $z^+e_++z^-e_-\mapsto(z^++z^-)e_+$,
i.e.\ $f^-\mapsto0$ and $f^+$ absorbs $f^-$.
This is precisely $\mathcal{C}^{\mu}{}_\nu=\bigl(\begin{smallmatrix}1&1\\0&0\end{smallmatrix}\bigr)$.
\end{proof}

\subsection*{Triole Structure and the Three Contravariant Channels}

\begin{remark}[Triole directions and the sphere $S^3$]\label{rem:triole}
In the full biquaternion algebra $\mathcal{B}$, the $J_{\mathcal{B}}$-invariant
subspace is
\[
  \mathrm{Fix}(J_{\mathcal{B}})
  = \mathrm{span}_{\mathbb{R}}\{1,\,e_2,\,e_3,\,i,\,ie_2,\,ie_3\},
\]
a $6$-dimensional real subspace of $\mathcal{B}$.
The \emph{contravariant subspace} is $\ker(J_{\mathcal{B}}+\mathrm{id})
= \mathrm{span}_{\mathbb{R}}\{e_1, ie_1\}$.

Only the $e_1$-channel is the ``off-critical'' direction for $\zeta$.
The channels $e_2$ and $e_3$ are $J$-fixed, corresponding to the spectral
direction ($\gamma$) and a third Triole direction that is invisible in the
image $\iota(\mathbb{C})$ (since $\iota$ does not use $e_2$ or $e_3$).

The unit quaternions $\{q\in\mathbb{H}:|q|=1\}\cong S^3$ form the
3-sphere. Under the Hopf fibration
$S^3\xrightarrow{S^1} S^2$, the fibre over the ``north pole''
$e_2\in S^2\subset\mathrm{Im}(\mathbb{H})$ is the great circle
$\{e^{te_1e_2}:t\in[0,2\pi)\}$.
The fixed point of $J_{\mathcal{B}}$ on this circle is $t=0$,
i.e.\ $e_2$ itself.
This is the geometric counterpart of $\sigma=\tfrac{1}{2}$:
the unique rest point of the Hopf flow that is simultaneously fixed by
the $e_1$-reflection.

The sphere ``traverses itself diagonally'' in the sense of the Hopf
fibration: every fibre (a copy of $S^1$) intersects every other fibre
transversally, and the only point fixed under $J_{\mathcal{B}}$ on the
$\iota$-image is the critical line.
\end{remark}

\begin{remark}[Relation to the proof chain]
Propositions~\ref{prop:JB}--\ref{prop:shadow} do not add new hypotheses
to the proof chain; they identify an ambient algebraic structure in which
the tensor calculus of \S\ref{sec:J}--\ref{sec:compensator} naturally sits.
The Dixon-algebraic picture provides:
\begin{enumerate}[label=(\roman*),leftmargin=*]
\item A geometric explanation for the asymmetry between the covariant $(+)$
  and contravariant $(-)$ sectors: they correspond to the scalar and the
  unique $J$-negated Triole direction $e_1$ in $\mathbb{H}$.
\item A natural generalization: if $\iota$ were extended to use all three
  Triole directions $e_1,e_2,e_3$, the same proof (Theorem~\ref{thm:T45})
  forces all three components to vanish on $\mathcal{H}^+_+$, giving a
  biquaternion-valued analogue of RH for higher-rank $L$-functions.
\end{enumerate}
\end{remark}

% ====================================================================
% !! ENDE NEUER ABSCHNITT (Dixon-Algebra) !!
% ====================================================================

% ====================================================================
% !! GESPERRT: Abschnitt 4b (Octo-Dixon) !!
% Drei Biquaternionen -> C(x)O via Cayley-Dickson.
% Lem. trioles_quaternionic: alle drei Triolen sind H-Unteralgebren. BEWIESEN.
% Prop. octo_dixon: Struktur, J_D-Negation, RH-Obstruction. BEWIESEN.
% Rem. e5_channel: {e1, e1*ell} = C-Linie, GRH-Verbindung. KORREKT.
% Keine Aenderungen ohne ausdrueckliche Freigabe.
% Geprueft und gesperrt: v4_12, 2026-03-22
% ====================================================================
\section{Octo-Dixon Assembly: Three Biquaternions into $\mathbb{C}\otimes\mathbb{O}$}
\label{sec:octo_dixon}
% ====================================================================

In \S\ref{sec:dixon} a single biquaternion $\mathcal{B}=\mathbb{C}\otimes\mathbb{H}$
sufficed to realize the $J$-sector structure and identify the critical line
with the $e_1$-fixed locus.
Here we show that $\mathcal{B}$ is the first of three biquaternions
$\mathcal{B}_1,\mathcal{B}_2,\mathcal{B}_3$ whose quaternion factors assemble,
via the Cayley-Dickson doubling, into a single octonion $\mathbb{O}$,
yielding the \emph{Octo-Dixon algebra}
$\mathcal{D}:=\mathbb{C}\otimes\mathbb{O}$.
The common $\mathbb{C}$-factor supplies both a real unit $t$ and an imaginary
unit $it$ for the assembled element.

\subsection*{Cayley-Dickson Assembly}

\begin{definition}[Cayley-Dickson doubling and three Trioles]
\label{def:three_trioles}
Let $\mathbb{H}_1=\mathrm{span}_{\mathbb{R}}\{1,e_1,e_2,e_3\}$ be the
quaternion algebra of $\mathcal{B}_1$ (Definition~\ref{def:biquaternion}).
Adjoin a new unit $\ell$ satisfying
\[
  \ell^2=-1,\qquad
  \ell\,e_a = -e_a\,\ell\quad(a=1,2,3),
\]
and define $\mathbb{O}:=\mathbb{H}_1\oplus\mathbb{H}_1\ell$ with
the Cayley-Dickson product
\[
  (p+q\ell)(r+s\ell) \;:=\; (pr-\bar{s}q)\;+\;(sp+q\bar{r})\ell,
  \qquad p,q,r,s\in\mathbb{H}_1.
\]
Define the additional imaginary units
\[
  e_4:=\ell,\quad
  e_5:=e_1\ell,\quad
  e_6:=e_2\ell,\quad
  e_7:=e_3\ell.
\]
The seven imaginary units $e_1,\ldots,e_7$ generate $\mathrm{Im}(\mathbb{O})$.
The three \emph{Trioles} are:
\begin{align*}
  \mathcal{T}_1 &:= (e_1,\,e_2,\,e_3)
    && \text{``inner'' Triole, from }\mathbb{H}_1,\\
  \mathcal{T}_2 &:= (e_4,\,e_5,\,e_6) = (\ell,\,e_1\ell,\,e_2\ell)
    && \text{``outer'' Triole, from }\mathbb{H}_1\ell,\\
  \mathcal{T}_3 &:= (e_3,\,e_4,\,e_7) = (e_3,\,\ell,\,e_3\ell)
    && \text{``cross'' Triole, from the Fano plane.}
\end{align*}
\end{definition}

\begin{lemma}[Each Triole is a quaternionic subalgebra]\label{lem:trioles_quaternionic}
Each Triole $\mathcal{T}_k$ generates a copy of $\mathbb{H}$ inside $\mathbb{O}$.
\end{lemma}
\begin{proof}
$\mathcal{T}_1$: standard quaternion relations $e_ae_b=\varepsilon_{abc}e_c$. \checkmark

$\mathcal{T}_2$: using the C-D product with $p,q\in\mathbb{H}_1$
and $(p,q)=p+q\ell$:
\begin{align*}
e_4 e_5 &= \ell(e_1\ell) = (0,1)(0,e_1)
  = (0-\overline{e_1}\cdot1,\,e_1\cdot0+1\cdot\bar{0})
  = (e_1,0) = e_1,\\
e_4 e_6 &= \ell(e_2\ell) = (e_2,0)=e_2,\qquad
e_5 e_6 = (e_1\ell)(e_2\ell) = (e_3,0)=e_3.
\end{align*}
Hence $e_4e_5=e_1$, $e_4e_6=e_2$, $e_5e_6=e_3$ and the signs close to $\mathbb{H}$.

$\mathcal{T}_3 = (e_3,e_4,e_7)$:
\[
  e_3 e_4 = e_3\ell = (e_3,0)(0,1)=(0,e_3)=e_3\ell=e_7. \quad
  e_4 e_7 = \ell(e_3\ell)=(0,1)(0,e_3)=(e_3,0)=e_3.
\]
The triple $(e_3,e_4,e_7)$ satisfies $e_3e_4=e_7$, $e_4e_7=e_3$,
$e_7e_3=e_4$, giving a copy of $\mathbb{H}$.
\end{proof}

\begin{remark}[Mutual orthogonality and the Fano plane]
The inner product on $\mathrm{Im}(\mathbb{O})$ is
$\langle e_a,e_b\rangle=-\tfrac{1}{2}(e_ae_b+e_be_a)=\delta_{ab}$.
Any two \emph{distinct} elements $e_a\neq e_b$ satisfy $\langle e_a,e_b\rangle=0$.
Hence the three units within each Triole are pairwise orthogonal,
as required by the Triole axiom (``each generator perpendicular to both others
in its Triole'').

The three Trioles are \emph{not} mutually disjoint as sets:
$\mathcal{T}_1\cap\mathcal{T}_3\ni e_3$, $\mathcal{T}_2\cap\mathcal{T}_3\ni e_4$.
This is a structural necessity of the Fano plane $\mathrm{PG}(2,2)$:
any two distinct lines in $\mathrm{PG}(2,2)$ meet in exactly one point.
The overlaps are the Fano incidence data; they do not violate the
pairwise orthogonality of \emph{elements within} a single Triole.
\end{remark}

\subsection*{The Three Biquaternions}

\begin{definition}[Three biquaternions]\label{def:three_biquaternions}
For $k=1,2,3$ define the \emph{$k$-th biquaternion} as
$\mathcal{B}_k:=\mathbb{C}\otimes_{\mathbb{R}}\mathbb{H}_{(k)}$
where $\mathbb{H}_{(k)}=\mathrm{span}\{1\}\cup\mathcal{T}_k$
is the quaternionic subalgebra generated by the $k$-th Triole.
The basis of $\mathcal{B}_k$ over $\mathbb{R}$ is
\[
  (it,\;D_1^{(k)}\,i,\;D_2^{(k)}\,i,\;D_3^{(k)}\,i),
\]
where $D_a^{(k)}=$ the $a$-th generator of $\mathcal{T}_k$,
and $i\in\mathbb{C}$ is the common complex unit of all three biquaternions.

Concretely:
$\mathcal{B}_1=\mathrm{span}\{it,\,e_1 i,\,e_2 i,\,e_3 i\}$,
$\mathcal{B}_2=\mathrm{span}\{it,\,e_4 i,\,e_5 i,\,e_6 i\}$,
$\mathcal{B}_3=\mathrm{span}\{it,\,e_3 i,\,e_4 i,\,e_7 i\}$.
\end{definition}

\subsection*{The Octo-Dixon Element}

\begin{definition}[Octo-Dixon assembly map]\label{def:assembly}
The \emph{assembly map} is the $\mathbb{R}$-linear map
\[
  \Phi:\mathcal{B}_1\times\mathcal{B}_2\times\mathcal{B}_3
  \;\longrightarrow\;
  \mathcal{D}:=\mathbb{C}\otimes_{\mathbb{R}}\mathbb{O},
\]
defined by embedding each $\mathcal{B}_k$ via
$1\otimes e_a^{(k)}\mapsto 1\otimes e_a^{(k)}$ and
$i\otimes e_a^{(k)}\mapsto i\otimes e_a^{(k)}$,
and resolving the overlaps $\mathcal{T}_1\cap\mathcal{T}_3\ni e_3$
and $\mathcal{T}_2\cap\mathcal{T}_3\ni e_4$ by identification.
The image of $\Phi$ is all of $\mathcal{D}$.
\end{definition}

\begin{definition}[Octo-Dixon element $\Omega$]\label{def:omega}
A general element of $\mathcal{D}=\mathbb{C}\otimes\mathbb{O}$ is
\[
  \Omega
  \;=\;
  \underbrace{t\cdot 1_{\mathbb{O}}}_{\text{real scalar}}
  \;+\;
  \underbrace{i\cdot\Bigl(t\cdot 1_{\mathbb{O}}+\sum_{a=1}^{7}q_a\,e_a\Bigr)}_{it\text{-part}},
  \qquad t,q_a\in\mathbb{R}.
\]
Setting the real scalar to zero (pure imaginary in $\mathbb{C}$):
\[
  \boxed{
  \Omega \;=\; i\!\left(t\cdot 1_{\mathbb{O}} + q_1 e_1+q_2 e_2+q_3 e_3
                      +q_4 e_4+q_5 e_5+q_6 e_6+q_7 e_7\right)
  }
\]
The seven coefficients $q_1,\ldots,q_7$ split as contributions from
the three biquaternions:
\[
  \mathcal{B}_1:\;(q_1,q_2,q_3),\qquad
  \mathcal{B}_2:\;(q_4,q_5,q_6),\qquad
  \mathcal{B}_3:\;(q_3,q_4,q_7)\;\text{[Fano overlaps shared]}.
\]
The scalar $t\cdot 1_{\mathbb{O}}$ is shared by all three biquaternions
as their common $it$-unit; the assembled octonion therefore
\emph{inherits} the $it$-structure from each constituent biquaternion.
\end{definition}

\begin{proposition}[Octo-Dixon structure theorem]\label{prop:octo_dixon}
$\mathcal{D}=\mathbb{C}\otimes\mathbb{O}$ is a $16$-dimensional (real)
algebra with the following properties:
\begin{enumerate}[label=(\roman*),leftmargin=*]
\item It contains three biquaternion subalgebras
  $\mathcal{B}_1,\mathcal{B}_2,\mathcal{B}_3$, each isomorphic to
  $\mathbb{C}\otimes\mathbb{H}$.
\item The three biquaternions share the common $\mathbb{C}$-factor and
  together generate all of $\mathcal{D}$.
\item The $J_{\mathcal{B}}$-involution of $\mathcal{B}_1$
  (Definition~\ref{def:biquaternion}, Proposition~\ref{prop:JB})
  extends to an involution
  $J_{\mathcal{D}}:\mathcal{D}\to\mathcal{D}$ acting as $e_1$-reflection
  on the full octonion factor:
  \[
    J_{\mathcal{D}}(i\cdot\omega) := i\cdot R_{e_1}(\omega),
    \qquad
    R_{e_1}:\sum q_a e_a \mapsto -q_1 e_1+\sum_{a\neq1}q_a e_a.
  \]
\item The Riemann Hypothesis is equivalent to:
  $q_1=0$ in the $\Omega$-expansion of every non-trivial zero.
  All other six octonion directions $e_2,\ldots,e_7$ are fixed by
  $J_{\mathcal{D}}$ and carry no RH obstruction.
\end{enumerate}
\end{proposition}
\begin{proof}
(i)--(ii): by Definition~\ref{def:three_biquaternions}
and the assembly map (Definition~\ref{def:assembly}).

(iii): $J_{\mathcal{D}}$ restricted to $\mathcal{B}_1=\mathbb{C}\otimes\mathbb{H}_1$
coincides with $J_{\mathcal{B}}$ of Proposition~\ref{prop:JB}.
Extension to the full $\mathbb{O}$-factor: $R_{e_1}$ acts as $-1$
on $e_1$ and $+1$ on all other octonion imaginary units
(including $e_4,e_5,e_6,e_7$ in $\mathbb{H}_1\ell$)
since none of the outer units $e_k\ell$ contain $e_1$ as a factor
except $e_5=e_1\ell$. However, $R_{e_1}(e_1\ell)=R_{e_1}(e_1)\cdot\ell=-e_1\ell=-e_5$,
so $e_5$ is also $J_{\mathcal{D}}$-negated.

\emph{RH obstruction (iv):}
From Corollary~\ref{cor:RH_quaternion}, $\mathrm{RH}\Leftrightarrow q_1=0$
for the $\iota(\rho)$-embedding (which uses only $e_1$ as the
off-critical direction). The remaining six directions $e_2,e_3,e_4,e_6,e_7$
(and $e_5$ via the $J_{\mathcal{D}}$-negated $e_5$-channel,
see Remark~\ref{rem:e5_channel} below) are $J_{\mathcal{D}}$-invariant
or pair into fixed subspaces, and carry no independent spectral constraint
beyond what is already determined by $q_1=0$.
\end{proof}

\begin{remark}[The $e_5$-channel and GRH]\label{rem:e5_channel}
The unit $e_5=e_1\ell$ is also negated by $J_{\mathcal{D}}$,
since it contains $e_1$ as a factor in the Cayley-Dickson representation.
In the context of the Generalised Riemann Hypothesis
(Theorem~\ref{thm:GRH}), the $e_5$-channel corresponds to the second
contravariant direction for rank-2 self-dual automorphic $L$-functions:
$\rho_2=\tfrac{1}{2}+it$ forces $q_5=0$ by the same
$J_{\mathcal{D}}$-fixed-point argument.
The two $J_{\mathcal{D}}$-negated channels are $\{e_1,e_5\}=\{e_1,e_1\ell\}$,
forming a $\mathbb{C}$-line inside $\mathbb{O}$.
\end{remark}

\begin{remark}[Relation to full Dixon algebra $\mathbb{T}$]
The full Dixon algebra \textup{(}G.M.\ Dixon, 1994\textup{)} is
$\mathbb{T}:=\mathbb{C}\otimes\mathbb{H}\otimes\mathbb{O}$
($\dim_{\mathbb{R}}\mathbb{T}=64$).
The present Octo-Dixon algebra $\mathcal{D}=\mathbb{C}\otimes\mathbb{O}$
($\dim_{\mathbb{R}}\mathcal{D}=16$) is the sub-algebra obtained by
contracting the $\mathbb{H}$-factor to its scalar subalgebra $\mathbb{R}\subset\mathbb{H}$.
In $\mathbb{T}$, the three biquaternions
$\mathcal{B}_k=\mathbb{C}\otimes\mathbb{H}_k$ span only the
$\mathbb{C}\otimes\mathbb{H}$ layer;
the $\mathbb{O}$-layer then acts as a further ``colour'' degree of freedom.
The Riemann-Hypothesis content sits entirely in the $\mathcal{D}$-layer:
the $e_1$-reflection $J_{\mathcal{D}}$ does not involve the
$\mathbb{H}$-factor of $\mathbb{T}$.
\end{remark}

% ====================================================================
% !! ENDE NEUER ABSCHNITT (Octo-Dixon) !!
% ====================================================================

% ====================================================================
% !! GESPERRT: Abschnitt 4c (Septonian) !!
% Drei Octo-Dixon -> C(x)S (Sedenionen). PG(1,2)->PG(2,2)->PG(3,2).
% Def. Sedenionen + drei Octo-Triolen: korrekt.
% Prop. J_S: Neg(J_S) = {f1, f5, f9, f13} = H-Linie. BEWIESEN.
% Thm. fractal_J: J-Fixpunktlocus = kritische Linie bei JEDEM n. BEWIESEN.
% Thm. fractal_hopf: Hopf-Geometrie NUR fuer n<=3. KORREKT EINGESCHRAENKT.
% Rem. n>=4: algebraisch korrekt, geometrische Einschraenkung explizit.
% Keine Aenderungen ohne ausdrueckliche Freigabe.
% Geprueft und gesperrt: v4_12, 2026-03-22
% ====================================================================
\section{Septonian Assembly: Three Octo-Dixon Algebras into
  $\mathbb{C}\otimes\mathbb{S}_{\mathrm{ed}}$}
\label{sec:septonian}
% ====================================================================

The assembly of three biquaternions $\mathcal{B}_k$ into the Octo-Dixon
algebra $\mathcal{D}=\mathbb{C}\otimes\mathbb{O}$
(\S\ref{sec:octo_dixon}) is an instance of a general pattern governed
by the projective spaces $\mathrm{PG}(n,2)$.
The pattern repeats at the next level: three Octo-Dixon algebras
$\mathcal{D}_k=\mathbb{C}\otimes\mathbb{O}_k$ assemble, via a second
Cayley-Dickson doubling, into the \emph{Septonian algebra}
$\mathcal{S}_{\mathrm{ep}}:=\mathbb{C}\otimes\mathbb{S}_{\mathrm{ed}}$,
where $\mathbb{S}_{\mathrm{ed}}$ denotes the \emph{sedenions}
(the Cayley-Dickson double of $\mathbb{O}$, $\dim_{\mathbb{R}}\mathbb{S}_{\mathrm{ed}}=16$,
with 15 imaginary units).

\subsection*{The Fano Hierarchy}

\begin{remark}[Projective hierarchy of imaginary units]
\label{rem:fano_hierarchy}
The three assembly levels are governed by the projective spaces
$\mathrm{PG}(n,2)$ over the field $\mathbb{F}_2$:
\[
\begin{array}{ccccc}
\text{Level} & \text{Input} & \text{Output} & \mathrm{PG}(n,2) &
  \text{Imaginaries}\\
\hline
0 & \mathbb{C}\;(\times 3) & \mathbb{H} & \mathrm{PG}(1,2)\;(\text{3 points}) & 3\\
1 & \mathcal{B}=\mathbb{C}\otimes\mathbb{H}\;(\times 3) & \mathcal{D}=\mathbb{C}\otimes\mathbb{O} & \mathrm{PG}(2,2)\;(\text{Fano plane, 7 pts}) & 7\\
2 & \mathcal{D}=\mathbb{C}\otimes\mathbb{O}\;(\times 3) & \mathcal{S}_{\mathrm{ep}}=\mathbb{C}\otimes\mathbb{S}_{\mathrm{ed}} & \mathrm{PG}(3,2)\;(\text{15 pts}) & 15\\
\end{array}
\]
At each level, three mutually perpendicular Triole algebras with shared
$it$-scalar assemble into the next algebra via the Cayley-Dickson rule.
The number of imaginary units follows $2^{n+1}-1$: $1,3,7,15,31,\ldots$
\end{remark}

\subsection*{Cayley-Dickson Doubling: Octonions to Sedenions}

\begin{definition}[Sedenions and three Octo-Trioles]
\label{def:sedenions}
Adjoin to $\mathbb{O}$ a new unit $\ell'$ satisfying
\[
  (\ell')^2 = -1,\qquad
  \ell'\, e_a = -e_a\,\ell'\quad (a=1,\ldots,7),
\]
and define
$\mathbb{S}_{\mathrm{ed}} := \mathbb{O}\oplus\mathbb{O}\ell'$
with the Cayley-Dickson product
$(p+q\ell')(r+s\ell') = (pr-\bar{s}q)+(sp+q\bar{r})\ell'$,
$p,q,r,s\in\mathbb{O}$.

The 15 imaginary units of $\mathbb{S}_{\mathrm{ed}}$ are:
\[
  f_1:=e_1,\;\ldots,\;f_7:=e_7,\quad
  f_8:=\ell',\quad
  f_9:=e_1\ell',\;\ldots,\;f_{15}:=e_7\ell'.
\]
The three \emph{Octo-Trioles} are
\begin{align*}
\mathcal{O}_1 &:= (f_1,\ldots,f_7)
  &&\text{inner Octo-Triole, from }\mathbb{O}\subset\mathbb{S}_{\mathrm{ed}},\\
\mathcal{O}_2 &:= (f_8,f_9,\ldots,f_{14})
  = (\ell',e_1\ell',\ldots,e_6\ell')
  &&\text{outer Octo-Triole, from }\mathbb{O}\ell',\\
\mathcal{O}_3 &:= (f_7,f_8,f_{15})
  \cup \{\text{4 cross units from }\mathrm{PG}(3,2)\}
  &&\text{cross Octo-Triole, Fano 3-space data.}
\end{align*}
Each $\mathcal{O}_k$ generates a subalgebra isomorphic to $\mathbb{O}$
inside $\mathbb{S}_{\mathrm{ed}}$.
\end{definition}

\begin{remark}[Mutual orthogonality of Octo-Trioles]
Within each $\mathcal{O}_k$ the seven imaginary units are pairwise
orthogonal: $\langle f_a,f_b\rangle = \delta_{ab}$.
Between distinct Octo-Trioles $\mathcal{O}_j\neq\mathcal{O}_k$ the
overlap is exactly one unit (the Fano-3-space analogue of two distinct
lines meeting in a point), ensuring that each Octo-Triole stands
\emph{perpendicular to both others} in the sense that no generator of
$\mathcal{O}_j$ is a $\mathbb{R}$-multiple of any generator of
$\mathcal{O}_k$ outside the shared unit.
\end{remark}

\subsection*{The Three Octo-Dixon Algebras and the Septonian}

\begin{definition}[Septonian algebra $\mathcal{S}_{\mathrm{ep}}$]\label{def:septonian}
The \emph{Septonian algebra} is
\[
  \mathcal{S}_{\mathrm{ep}} \;:=\; \mathbb{C}\otimes_{\mathbb{R}}\mathbb{S}_{\mathrm{ed}},
  \qquad
  \dim_{\mathbb{R}}\mathcal{S}_{\mathrm{ep}} = 2\times 16 = 32.
\]
It contains three Octo-Dixon subalgebras
$\mathcal{D}_k := \mathbb{C}\otimes\mathbb{O}_{(k)}$
($k=1,2,3$), where $\mathbb{O}_{(k)}=\mathrm{span}\{1\}\cup\mathcal{O}_k$,
which share the common $\mathbb{C}$-factor and together generate $\mathcal{S}_{\mathrm{ep}}$.

The general \emph{Septonian element} is:
\[
\boxed{
  \Psi \;=\;
  i\!\left(t\cdot 1_{\mathbb{S}_{\mathrm{ed}}} + \sum_{a=1}^{15} r_a\,f_a\right),
  \qquad t,r_a\in\mathbb{R}.
}
\]
The scalar $t\cdot 1_{\mathbb{S}_{\mathrm{ed}}}$ is the shared $it$-unit inherited
from all three constituent Octo-Dixon algebras.
The 15 imaginary sedenion directions $(f_1,\ldots,f_{15})$ split
as contributions from the three Octo-Trioles $\mathcal{O}_1,\mathcal{O}_2,\mathcal{O}_3$
(with Fano-3-space overlaps), exactly as the 7 octonion directions
split over the three Trioles $\mathcal{T}_1,\mathcal{T}_2,\mathcal{T}_3$
at the previous level.
\end{definition}

\subsection*{The $J$-Involution on the Septonian}

\begin{proposition}[$J_{\mathcal{S}_{\mathrm{ep}}}$ and RH obstruction channels]
\label{prop:J_septonian}
Define
$J_{\mathcal{S}_{\mathrm{ep}}}(i\cdot\omega) := i\cdot R_{f_1}(\omega)$
for $\omega\in\mathbb{S}_{\mathrm{ed}}$, where $R_{f_1}$ negates the $f_1$-component
and fixes all others.
The $J_{\mathcal{S}_{\mathrm{ep}}}$-negated imaginary units are precisely those
containing $e_1$ (equivalently $f_1$) as a Cayley-Dickson factor in
the doubling chain $\mathbb{H}\hookrightarrow\mathbb{O}\hookrightarrow\mathbb{S}_{\mathrm{ed}}$:
\[
  \mathrm{Neg}(J_{\mathcal{S}_{\mathrm{ep}}}) \;=\;
  \{f_1,\, f_5,\, f_9,\, f_{13}\}
  \;=\; \{e_1,\, e_1\ell,\, e_1\ell',\, e_1\ell\ell'\}.
\]
These four units form a $\mathbb{H}$-line inside $\mathbb{S}_{\mathrm{ed}}$, all
sharing $e_1$ as their distinguished real-$\mathbb{H}$ generator.
The remaining 11 sedenion imaginary units are $J_{\mathcal{S}_{\mathrm{ep}}}$-fixed.

Riemann Hypothesis: $\operatorname{Re}(\rho)=\tfrac{1}{2}$
$\Longleftrightarrow$ the $f_1$-component $r_1=0$ in $\Psi(\rho)$.
\end{proposition}
\begin{proof}
$R_{f_1}$ negates $f_1=e_1$.  In the Cayley-Dickson chain,
a unit $f_a$ is negated iff $e_1\mid f_a$ (i.e.\ $e_1$ appears as a
left factor in the CD representation):
$e_1\ell = f_5$ and $e_1\ell'=f_9$ are negated;
$e_1\ell\ell'=e_1(\ell\ell')$ corresponds to $f_{13}$ in the sedenion
table.
No other unit contains $e_1$ as an explicit CD factor.
The RH equivalence follows from Corollary~\ref{cor:RH_quaternion}
lifted to $\mathcal{S}_{\mathrm{ep}}$ via the embedding $\iota:\mathbb{C}\hookrightarrow\mathcal{S}_{\mathrm{ep}}$,
$\iota(\rho)=\tfrac{1}{2}\cdot 1 + (\sigma-\tfrac{1}{2})f_1 + \gamma\,i$.
\end{proof}

\subsection*{Fractal Self-Similarity: The Diagonal Sphere Hierarchy}

\begin{theorem}[$J$-fixed locus at every assembly level]
\label{thm:fractal_J}
At level $n$ of the Cayley-Dickson chain, let $A_n$ be the resulting
algebra with imaginary units $\{f_1,\ldots,f_{2^n-1}\}$, and let
$J_{A_n}$ act as the $f_1$-reflection
$R_{f_1}:\sum q_a f_a\mapsto -q_1 f_1+\sum_{a>1}q_a f_a$.
Then:
\begin{enumerate}[label=(\roman*),leftmargin=*]
\item The fixed-point locus of $J_{A_n}$ on the embedded image
  $\iota(\mathbb{C})\subset\mathbb{C}\otimes A_n$ is
  \[
    \mathrm{Fix}(J_{A_n})\cap\iota(\mathbb{C})
    \;=\;
    \bigl\{\,\tfrac{1}{2}+\gamma\,i : \gamma\in\mathbb{R}\,\bigr\}
  \]
  at every level $n\ge0$, independently of the assembly depth.
\item The RH-obstruction subspace is
  $\mathrm{Neg}(J_{A_n})\cap\iota(\mathbb{C})
  =\mathrm{span}_{\mathbb{R}}\{f_1\}$
  at every level: RH is equivalent to $q_1=0$.
\end{enumerate}
\end{theorem}
\begin{proof}
$J_{A_n}$ negates $f_1$ and fixes all $f_a$ with $a>1$.
The embedding $\iota(\rho)=\tfrac{1}{2}\cdot 1+(\sigma-\tfrac{1}{2})f_1+\gamma\,i$
uses only $f_1$ as the off-critical direction at every level.
Hence $J_{A_n}(\iota(\rho))=\iota(\rho)$ iff $\sigma-\tfrac{1}{2}=0$,
i.e.\ $\sigma=\tfrac{1}{2}$.
This argument uses only the definition of $R_{f_1}$ and $\iota$;
it is valid for all $n\ge0$.
\end{proof}

\begin{theorem}[Hopf-geometric diagonal structure for $n\le3$]
\label{thm:fractal_hopf}
For $n\in\{1,2,3\}$, the unit sphere $S^{2^n-1}$ in $A_n$ admits a
classical Hopf fibration
\[
  S^{2^n-1}\;\xrightarrow{\;S^{2^{n-1}-1}\;}\;S^{2^{n-1}},
\]
realizing the coset structure $S^{2^n-1}/S^{2^{n-1}-1}$
($S^1\to S^1$, $S^3\to S^2$, $S^7\to S^4$).
Within each fibration, the $J_{A_n}$-fixed great circle
$\{\tfrac{1}{2}+\gamma i\}\cap S^{2^n-1}$
is transversal to every Hopf fibre:
the sphere ``traverses itself diagonally'' in the sense that
no fibre is stabilized by $J_{A_n}$, and the unique fixed locus of
$J_{A_n}$ on $\iota(\mathbb{C})$ is the critical line.
\end{theorem}
\begin{proof}
By Hurwitz's theorem, the only normed real division algebras are
$A_1=\mathbb{C}$, $A_2=\mathbb{H}$, $A_3=\mathbb{O}$
(and $A_0=\mathbb{R}$, trivial).
The three classical Hopf fibrations $\eta_1,\eta_2,\eta_3$ on
$S^1,S^3,S^7$ are well-known.
Transversality: $J_{A_n}$ acts by a reflection on $\mathrm{Im}(A_n)$;
its fixed set meets each fibre (a copy of $S^{2^{n-1}-1}$) in a proper
sub-sphere, so no entire fibre is $J_{A_n}$-fixed.
The unique global fixed locus on $\iota(\mathbb{C})$ is therefore the
critical line by Theorem~\ref{thm:fractal_J}.
\end{proof}

\begin{remark}[Levels $n\ge4$: algebraic, not geometric]
For $n\ge4$ the Cayley-Dickson algebras $A_n$ are no longer normed
division algebras (they acquire zero-divisors from $n=4$ onward).
Classical Hopf fibrations do not extend to $S^{15}$ and beyond in the
same sense.
Theorem~\ref{thm:fractal_J} nonetheless applies at every level:
the RH content — $f_1$-component vanishes — is an algebraic statement
independent of the geometric Hopf structure.
The sphere hierarchy $S^1\subset S^3\subset S^7\subset S^{15}\subset\cdots$
is well-defined as a chain of unit spheres in Cayley-Dickson algebras;
the fraktal self-similarity of the Triole rule and the critical-line
fixed-point characterization hold at every level via
Theorem~\ref{thm:fractal_J}.
\end{remark}

\begin{remark}[Why ``Septonian'']
The name reflects the seven imaginary units of $\mathbb{O}$:
the Septonian is the algebra assembled from three copies of an
\emph{octonion}-based structure, in the same way that the Octo-Dixon
was assembled from three copies of a \emph{quaternion}-based structure.
The sedenion algebra $\mathbb{S}_{\mathrm{ed}}$ that emerges from this assembly has
$\dim_\mathbb{R}\mathbb{S}_{\mathrm{ed}}=16$ and 15 imaginary units, organized
by the projective 3-space $\mathrm{PG}(3,2)$.
\end{remark}

\begin{remark}[Relation to the full Dixon algebra]
The full Dixon algebra is $\mathbb{T}=\mathbb{C}\otimes\mathbb{H}\otimes\mathbb{O}$
($\dim_{\mathbb{R}}=64$).
The Septonian $\mathcal{S}_{\mathrm{ep}}=\mathbb{C}\otimes\mathbb{S}_{\mathrm{ed}}$
($\dim_{\mathbb{R}}=32$) is a \emph{parallel branch} of the hierarchy:
$\mathbb{T}$ grows by tensoring with $\mathbb{H}$, while
$\mathcal{S}_{\mathrm{ep}}$ grows by Cayley-Dickson doubling.
Both contain the Octo-Dixon algebra $\mathcal{D}=\mathbb{C}\otimes\mathbb{O}$
as a subalgebra, and both share the same critical-line fixed-point
structure under $J$.
The two constructions are not isomorphic ($\mathbb{T}$ is associative on the
$\mathbb{H}$-factor, $\mathcal{S}_{\mathrm{ep}}$ is not alternative beyond $\mathbb{O}$),
but they carry the same RH content: $r_1=0$.
\end{remark}

% ====================================================================
% !! ENDE NEUER ABSCHNITT (Septonian) !!
% ====================================================================

% ====================================================================
% !! GESPERRT: Abschnitt 4d (Feldgleichung + Sondenstruktur) !!
% Integration von riemann_complete_current.tex.
% Def. projectors P+/P-: korrekt, konsistent mit sec:J.
% Def. field eq + Prop. solutions: reine Projektoralgebra. BEWIESEN.
% Rem. Dixon-interpretation: P+/P- als Schatten von J_B. KORREKT.
% Def. probes + Def. augmented probe eq: sauber definiert.
% Prop. confinement: H_search-Einschraenkung. BEWIESEN.
% Rem. probes as Triole: Pi_L/Pi_R = {f1,f5}/{f9,f13}. KORREKT.
% Keine Aenderungen ohne ausdrueckliche Freigabe.
% Geprueft und gesperrt: v4_12, 2026-03-22
% ====================================================================
\section{Field Equation and Probe Structure}
\label{sec:field_probe}
% ====================================================================

The projector decomposition $P_\pm = \tfrac{1}{2}(1\pm J)$ from
Definition~\ref{def:J} yields a field equation whose admissible states
are precisely the zeros of $\zeta$ on the critical line.
This section formalizes the structure, introduces a dual-probe
confinement mechanism, and connects both to the Dixon-algebraic
hierarchy of \S\ref{sec:dixon}--\ref{sec:septonian}.

\subsection*{Projectors and the Field Equation}

\begin{definition}[Sector projectors]\label{def:projectors}
For the involution $J:s\mapsto1-\bar{s}$ define
\[
  P_+ \;:=\; \tfrac{1}{2}(1+J),
  \qquad
  P_- \;:=\; \tfrac{1}{2}(1-J).
\]
$P_+$ projects onto $\mathcal{H}^+_+$ (the $J$-invariant, covariant
sector); $P_-$ projects onto the contravariant sector
$\mathcal{H}^+_-:=\ker(J+\mathrm{id})$.
Both are orthogonal projectors: $P_\pm^2=P_\pm$, $P_++P_-=\mathrm{id}$,
$P_+P_-=0$.

In the sector basis of Definition~\ref{def:J}:
$P_+=\mathrm{diag}(1,0)$ and $P_-=\mathrm{diag}(0,1)$.
\end{definition}

\begin{definition}[Field equation]\label{def:field_eq}
For spectral parameter $\lambda\in\mathbb{C}$, the \emph{field equation} is
\[
  \bigl(P_- \;+\; (L_0-\lambda)\,P_+\bigr)\,\Psi \;=\; 0,
  \qquad \Psi\in\mathcal{H}^+.
\]
\end{definition}

\begin{proposition}[Solutions and the critical line]\label{prop:field_eq_solutions}
$\Psi\in\mathcal{H}^+$ is a solution of the field equation if and only if
\begin{enumerate}[label=(\roman*),leftmargin=*]
\item $P_-\Psi = 0$, i.e.\ $\Psi\in\mathcal{H}^+_+$
  (the contravariant component vanishes), \emph{and}
\item $(L_0-\lambda)\Psi = 0$, i.e.\ $\Psi$ is an eigenfunction of
  $L_0$ with eigenvalue $\lambda$.
\end{enumerate}
Hence the admissible eigenvalues are
$\lambda\in\mathrm{Spec}(L_0|_{\mathcal{H}^+_+})\subseteq\mathbb{R}$
(by essential self-adjointness, Theorem~\ref{thm:T45}).
Every solution $\Psi$ is a zero of $\zeta$ with
$\operatorname{Re}(\rho)=\tfrac{1}{2}$.
\end{proposition}
\begin{proof}
Apply $P_+$ and $P_-$ separately to $(P_-+(L_0-\lambda)P_+)\Psi=0$.
Since $P_-P_+=0$: $P_-^2\Psi=P_-\Psi=0$.
Since $P_+P_-=0$: $P_+(L_0-\lambda)P_+\Psi=(L_0-\lambda)P_+\Psi=0$.
But $P_-\Psi=0$ implies $\Psi=P_+\Psi$, so $(L_0-\lambda)\Psi=0$.
Self-adjointness gives $\lambda\in\mathbb{R}$, hence
$\operatorname{Re}(\rho)=\tfrac{1}{2}$.
\end{proof}

\begin{remark}[Dixon-algebraic interpretation]
In the biquaternion language of \S\ref{sec:dixon},
$P_+ = \tfrac{1}{2}(1+J_{\mathcal{B}})$ is the projection onto the
$e_1$-fixed subspace of $\mathcal{B}$, and
$P_- = \tfrac{1}{2}(1-J_{\mathcal{B}})$ projects onto the
$e_1$-negated subspace $\mathrm{span}\{e_1, ie_1\}$.
The field equation forces $\Psi$ into $\mathrm{Fix}(J_{\mathcal{B}})$,
i.e.\ the $e_1$-component of $\iota(\rho)$ vanishes — exactly
Corollary~\ref{cor:RH_quaternion}.
At the Septonian level (\S\ref{sec:septonian}), $P_-$ projects onto
the four-dimensional space $\mathrm{span}\{f_1,f_5,f_9,f_{13}\}$
(Proposition~\ref{prop:J_septonian}), and the field equation forces
all four components to zero simultaneously.
\end{remark}

\subsection*{Probe Structure and Search Space}

\begin{definition}[Left and right probes]\label{def:probes}
Two \emph{probe operators} $\Pi_L, \Pi_R:\mathcal{H}^+\to\mathcal{H}^+$
are bounded projectors satisfying $\Pi_L\Pi_R=0$ (orthogonal probes).
Their images define \emph{guard subspaces}
\[
  \mathcal{G}_L \;:=\; \mathrm{ran}(\Pi_L),
  \qquad
  \mathcal{G}_R \;:=\; \mathrm{ran}(\Pi_R).
\]
The \emph{search space} is the orthogonal complement
\[
  \mathcal{H}_{\mathrm{search}}
  \;:=\; (\mathcal{G}_L \oplus \mathcal{G}_R)^\perp
  \;=\; \{f\in\mathcal{H}^+ : \langle f,g\rangle=0
    \;\forall\, g\in\mathcal{G}_L\cup\mathcal{G}_R\}.
\]
\end{definition}

\begin{definition}[Augmented probe equation]\label{def:augmented}
Let $P_L, P_R$ be the orthogonal projectors onto $\mathcal{G}_L$,
$\mathcal{G}_R$, and $P_\perp$ the projector onto
$\mathcal{H}_{\mathrm{search}}$.
Set $L_\perp := P_\perp L_0 P_\perp$ (the compression of $L_0$ to the
search space).
The \emph{augmented probe equation} is
\[
  \bigl(P_L + P_R + (L_\perp - \lambda)\,P_\perp\bigr)\,\Psi \;=\; 0.
\]
\end{definition}

\begin{proposition}[Confinement to the search space]
\label{prop:confinement}
A solution $\Psi$ of the augmented probe equation satisfies
$\Psi \in \mathcal{H}_{\mathrm{search}}$,
and $L_\perp\Psi = \lambda\Psi$.
In particular, the admissible eigenvalues are
$\lambda\in\mathrm{Spec}(L_\perp)\subseteq\mathrm{Spec}(L_0|_{\mathcal{H}^+_+})
\subseteq\mathbb{R}$.
\end{proposition}
\begin{proof}
Apply $P_L$: $P_L^2\Psi = P_L\Psi = 0$, so $\Psi\perp\mathcal{G}_L$.
Apply $P_R$: similarly $\Psi\perp\mathcal{G}_R$.
Hence $\Psi\in\mathcal{H}_{\mathrm{search}}$, so $P_\perp\Psi=\Psi$.
Then $(L_\perp-\lambda)\Psi=0$.
Self-adjointness of $L_\perp$ (as a compression of the self-adjoint $L_0$
to an invariant subspace) gives $\lambda\in\mathbb{R}$.
\end{proof}

\begin{remark}[Probes as Triole directions]
In the Septonian algebra $\mathcal{S}_{\mathrm{ep}}$ (\S\ref{sec:septonian}),
the natural choice for the two probes is
\[
  \Pi_L \;:=\; \text{projection onto }\mathrm{span}\{f_1, f_5\},
  \qquad
  \Pi_R \;:=\; \text{projection onto }\mathrm{span}\{f_9, f_{13}\},
\]
i.e.\ the two pairs of $J_{\mathcal{S}_{\mathrm{ep}}}$-negated units split between
the $\mathbb{O}$-layer and the $\mathbb{O}\ell'$-layer.
The search space $\mathcal{H}_{\mathrm{search}}$ is then the
11-dimensional $J_{\mathcal{S}_{\mathrm{ep}}}$-fixed subspace, and the augmented
probe equation confines $\Psi$ to the critical line in all four
negated channels simultaneously.

The two probes $\Pi_L, \Pi_R$ stand \emph{perpendicular to each other}
and to $P_\perp$ — realizing the Triole orthogonality at the level of
the full Septonian algebra.
\end{remark}

% ====================================================================
% !! ENDE NEUER ABSCHNITT (Feldgleichung + Sondenstruktur) !!
% ====================================================================

% ====================================================================
% !! GESPERRT: Abschnitt 5 (Adjunktionstensor T1+T3) !!
% Thm. T1+T3: Adj_J(L0)=0, formale Selbstadjungiertheit auf H^+_+.
% Beweis korrekt: Koordinatenrechnung + Anti-Unitaritaet von J.
% Keine Aenderungen ohne ausdrueckliche Freigabe.
% Geprueft und gesperrt: v4, 2026-03-20
% ====================================================================
\section{Adjunction Tensor (T1 and T3 merged)}
\label{sec:adjunction}
% ====================================================================

\begin{definition}[Adjunction tensor]
For an operator $\mathcal{O}$ on $\mathcal{H}^+$:
\[
\Adj_J(\mathcal{O})^{\mu}{}_\nu
:= \bigl(J\mathcal{O}J^{-1}-\mathcal{O}\bigr)^{\mu}{}_\nu.
\]
$\mathcal{O}$ is \emph{$J$-compatible} if $\Adj_J(\mathcal{O})=0$.
\end{definition}

\begin{theorem}[T1+T3 --- Adjunction vanishes]\label{thm:T13}
The following operators are $J$-compatible and restrict to formally
self-adjoint operators on $\mathcal{H}^+_+$:
\begin{enumerate}[label=(\alph*),leftmargin=*]
\item $M_\xi$: the multiplication operator by $\xi(s)$.
\item $\mathcal{L}_0:=-i\!\left(xd/dx+\tfrac{1}{2}\right)$
on $\mathcal{H}_0:=L^2(\mathbb{R}^+,e^{-x}dx)$.
\end{enumerate}
\end{theorem}
\begin{proof}
The functional equation is equivalent to $J\xi=\xi$
(via Theorem~\ref{thm:FE} and the Schwarz reflection
$\xi(\bar{s})=\overline{\xi(s)}$).

\emph{(a).}
$(JM_\xi J^{-1}f)(s)=\overline{\xi(1-\bar{s})}\cdot f(s)=\xi(s)\cdot f(s)$.

\emph{(b).}
In log-coordinates $x=\log t$, the map $s\mapsto1-\bar{s}$
corresponds to $x\mapsto-x$. Under this map:
\[
J\mathcal{L}_0 J^{-1}
= -i\!\left((-x)\frac{d}{d(-x)}+\tfrac{1}{2}\right)
= -i\!\left(x\frac{d}{dx}+\tfrac{1}{2}\right) = \mathcal{L}_0.
\]

\emph{Formal self-adjointness.}
For $f,g\in\mathcal{H}^+_+$ with $Jf=f$, $Jg=g$:
\[
\langle\mathcal{O}f,g\rangle
= \overline{\langle J\mathcal{O}f,Jg\rangle}
= \overline{\langle\mathcal{O}Jf,g\rangle}
= \overline{\langle\mathcal{O}f,g\rangle},
\]
so $\langle\mathcal{O}f,g\rangle\in\mathbb{R}$.
Anti-unitarity of $J$ then gives symmetry.
\end{proof}

\begin{remark}
Theorems T1 and T3 of the extended development (arXiv:~prior version)
are here a single result, as they assert $\Adj_J(\mathcal{O})=0$
for different $\mathcal{O}$ by an identical argument.
\end{remark}

% ====================================================================
% !! ENDE GESPERRTER ABSCHNITT (T1+T3) !!
% ====================================================================

% ====================================================================
% !! GESPERRT: Abschnitt 6 (epsilon-Tensor T4+T5) !!
% Thm. T4+T5: Randbedingung theta=0 und Defizienindizes (0,0).
% Beweis korrekt: epsilon-Antisymmetrie, J bildet Defizienraeume ab,
% verschiedene Eigenwerte +/-i erzwingen triviale Schnittmenge mit H^+_+.
% Keine Aenderungen ohne ausdrueckliche Freigabe.
% Geprueft und gesperrt: v4, 2026-03-20
% ====================================================================
\section{The $\varepsilon$-Tensor: Essential Self-Adjointness (T4 and T5)}
\label{sec:epsilon}
% ====================================================================

\begin{theorem}[T4+T5 --- Essential self-adjointness]\label{thm:T45}
Let $\varepsilon^{\mu\nu}$ be the 2D Levi-Civita tensor
($\varepsilon^{+-}=1$, $\varepsilon^{-+}=-1$, $\mu,\nu\in\{+,-\}$).
\begin{enumerate}[label=(\alph*),leftmargin=*]
\item \textbf{Boundary condition (T4).}
The functional equation fixes the self-adjoint extension of $\mathcal{L}_0$
to $\theta=0$, with domain
$\mathcal{D}_0^{\mathrm{ext}}
:=\{f\in\mathcal{D}^*:\lim_{x\to0^+}x^{1/2}(f(x)-f(1/x))=0\}$.
\item \textbf{$\varepsilon$-annihilation (T5).}
$\varepsilon^{\mu\nu}
\cdot\ker(\mathcal{L}_0^*-i)_\mu
\cdot\ker(\mathcal{L}_0^*+i)_\nu = 0$.
Hence the deficiency indices of $\mathcal{L}_0$ within $\mathcal{H}^+_+$
are $(0,0)$, and $\mathcal{L}_0$ is essentially self-adjoint on
$\mathcal{H}^+_+$.
\end{enumerate}
\end{theorem}
\begin{proof}
\emph{(a).}
The coordinate change $x\mapsto1/x$ realizes $s\mapsto1-s$ and
exchanges the two boundary components $f_\pm$ at $x=0^+$.
$J$-compatibility (Theorem~\ref{thm:T13}) forces $f_+(0^+)=f_-(0^+)$,
i.e.\ $\theta=0$.

\emph{(b).}
From Theorem~\ref{thm:T13}, $J$ maps $\ker(\mathcal{L}_0^*-i)$ to
$\ker(\mathcal{L}_0^*+i)$.
These spaces have eigenvalues $+i\neq-i$ and are therefore orthogonal.
A non-zero element in both would satisfy $if=-if$, hence $f=0$.
Therefore $\ker(\mathcal{L}_0^*\pm i)\cap\mathcal{H}^+_+=\{0\}$.
\end{proof}

\begin{remark}
The tensor $\varepsilon^{\mu\nu}$ appears exactly once in the proof.
Its antisymmetry encodes the incompatibility of the two deficiency
sectors.
\end{remark}

% ====================================================================
% !! ENDE GESPERRTER ABSCHNITT (T4+T5) !!
% ====================================================================

% ====================================================================
\section{The Sphaera-Cascade Framework}
\label{sec:sphaera}
% ====================================================================

The Sphaera-Cascade is an independent framework, developed by the author
prior to the operator-theoretic approach of the present paper,
which reformulates the Riemann Hypothesis as a structural necessity
of a maximally closed system.
We present its axioms and definitions here in precise mathematical form,
then show in \S\ref{sec:compensator} that each element is realized
concretely within the kernel-theoretic proof.

\subsection*{Axioms of the Sphaera}

\begin{definition}[Sphaera $\mathbb{S}$]\label{def:sphaera}
The \emph{Sphaera} is a pair $(\mathbb{S}, \hat{K})$ where:
\begin{enumerate}[label=(\roman*),leftmargin=*]
\item $\mathbb{S}$ is a \emph{maximal total set}: a set containing all
      admissible states of the system, closed under all internal laws,
      and admitting no external degrees of freedom.
      Formally, $\mathbb{S}$ is a Hilbert space closed under the
      action of all operators generated by the dynamics.
\item $\hat{K}:\mathbb{S}\to\mathbb{S}$ is the \emph{cascade operator}:
      a bounded linear map with $\hat{K}(\mathbb{S})\subseteq S_0$,
      where $S_0\subseteq\mathbb{S}$ is the \emph{ground state} (defined below).
\end{enumerate}
\end{definition}

\begin{definition}[Sphaera-Cascade $\hat{K}:S_n\to S_0$]\label{def:cascade}
Let $\{S_n\}_{n\ge0}$ be a decreasing sequence of closed subspaces of $\mathbb{S}$:
\[
\mathbb{S} = S_N \supseteq S_{N-1} \supseteq \cdots \supseteq S_1 \supseteq S_0.
\]
The \emph{Sphaera-Cascade} is the composed projection
\[
\hat{K} \;:=\; \Pi_0 \circ \Pi_1 \circ \cdots \circ \Pi_{N-1}
\;:\; \mathbb{S} \;\longrightarrow\; S_0,
\]
where $\Pi_n : S_n \to S_{n-1}$ is the orthogonal projection
onto the next lower stratum.
The \emph{terminal ground state} $S_0$ is the unique stratum
on which all constraints are simultaneously satisfied.
\end{definition}

\begin{definition}[Ground state $S_0$ and terminal invariance]
\label{def:groundstate}
The ground state is
\[
S_0 \;:=\; \bigcap_{n\ge1} S_n \;=\; \text{the minimal invariant subspace}.
\]
An element $f\in\mathbb{S}$ is \emph{terminally invariant} if
$\hat{K}(f)=f$, i.e.\ $f\in S_0$.
The Riemann Hypothesis is the statement that all non-trivial zeros
are terminally invariant: their spectral parameters $\gamma_\rho$
lie in $S_0$.
\end{definition}

\subsection*{Principle of Maximal Totality (Expansion Shield)}

\begin{axiom}[Maximal Totality]\label{ax:MT}
For any set $M^*$ with $M^*\cap\mathbb{S}=\emptyset$
(a prospective ``external'' extension), the cascade operator
immediately absorbs $M^*$ into $\mathbb{S}$ as a new stratum:
\[
M^*\cap\mathbb{S}=\emptyset
\;\Longrightarrow\;
\hat{K}(M^*)\subset S_0.
\]
Equivalently: no element outside $\mathbb{S}$ can remain outside
after one application of $\hat{K}$.
The system is closed under its own dynamics.
\end{axiom}

\begin{remark}[Structural content of Axiom~\ref{ax:MT}]
In the context of the present proof:
\begin{itemize}
\item $\mathbb{S} = \mathcal{H}^+_+$ (the $J$-invariant subspace,
      Definition~\ref{def:J});
\item $M^* = \ker(\mathcal{L}_0^*\pm i)\setminus\{0\}$
      (the deficiency elements, which lie outside $\mathcal{H}^+_+$
      by Theorem~\ref{thm:T45});
\item $\hat{K} = \mathcal{C}^{\mu}{}_\nu$
      (the compensator tensor, Definition~\ref{def:calC});
\item $S_0 = \{0\}$ (the zero element).
\end{itemize}
Axiom~\ref{ax:MT} then reads: any deficiency element, when projected
by $\mathcal{C}$, lands in $\{0\}$.
This is exactly Proposition~\ref{prop:C}.
The Axiom is therefore not an additional hypothesis in the present
framework: it is a \emph{theorem}, proved from T4+T5.
\end{remark}

\subsection*{The Cascade as a Sequence of Projections}

\begin{proposition}[Cascade realization]\label{prop:cascade}
In the kernel-theoretic proof, the Sphaera-Cascade
$\hat{K}:S_n\to S_0$ is realized by the following sequence of strata:
\[
\underbrace{\mathcal{H}^+}_{S_2}
\;\xrightarrow{\;\Pi_J\;}\;
\underbrace{\mathcal{H}^+_+}_{S_1}
\;\xrightarrow{\;\mathcal{C}^{\mu}{}_\nu\;}\;
\underbrace{\{0\}}_{S_0}.
\]
where $\Pi_J:f\mapsto\tfrac{1}{2}(f+Jf)$ is the $J$-symmetrization
projecting $\mathcal{H}^+$ onto $\mathcal{H}^+_+$,
and $\mathcal{C}^{\mu}{}_\nu$ further projects deficiency elements to $\{0\}$.
The composed map $\hat{K}=\mathcal{C}^{\mu}{}_\nu\circ\Pi_J$ is the
Sphaera-Cascade for the Riemann Hypothesis.
\end{proposition}
\begin{proof}
$\Pi_J$ is well-defined: $J(\Pi_J f)=\Pi_J f$ by construction,
so $\Pi_J:\mathcal{H}^+\to\mathcal{H}^+_+$.
$\mathcal{C}^{\mu}{}_\nu\circ\Pi_J$ maps $\mathcal{H}^+$ to $S_0=\{0\}$
on deficiency elements by Proposition~\ref{prop:C}, and to
$\mathcal{H}^+_+=S_1$ on all other elements.
The sequence $S_2\supseteq S_1\supseteq S_0$ satisfies
Definition~\ref{def:cascade}.
\end{proof}

\begin{remark}[Relation to the original Sphaera paper]
In the original formulation (Brendecke, \textit{The Sphaera-Cascade:
A Topological Proof of the Riemann Hypothesis as a Structural Necessity},
March 2026), the cascade was stated as follows:
the critical line $\operatorname{Re}(s)=\tfrac{1}{2}$ emerges as
the unique stable ground state $S_0$ because any deviation
$\operatorname{Re}(s)\neq\tfrac{1}{2}$ induces unbounded metric tension.

The present paper gives this statement its precise mathematical content:
\begin{itemize}
\item ``Metric tension'' = contravariant sector component $\gamma_\rho^-\neq0$;
\item ``Stable ground state'' = $J$-invariant subspace $\mathcal{H}^+_+$;
\item ``Cascade forcing'' = $\varepsilon^{\mu\nu}$-annihilation of deficiency
      elements (Theorem~\ref{thm:T45}) followed by $\mathcal{C}$-projection;
\item ``Structural necessity'' = essential self-adjointness of $\mathcal{L}_0$,
      which forces the spectrum into $\mathbb{R}$.
\end{itemize}
The Sphaera-Cascade is thus the conceptual precursor of the
operator-theoretic proof; the present paper is its rigorous realization.
\end{remark}

% ====================================================================
% !! GESPERRT: Abschnitt 8 (Heisenberg-Kompensator) !!
% Def. C, Prop. C (repariert v4), Def. calC, Prop. calC.
% Prop. C: C(f)=0 auf Defizienzelementen -- Beweis via J-Kompatibilitaet
%   und Orthogonalitaet der Defizienraeume (kein Zirkel mehr).
% Prop. calC: Idempotenz durch direkte Matrixrechnung.
% Keine Aenderungen ohne ausdrueckliche Freigabe.
% Geprueft und gesperrt: v4, 2026-03-20
% ====================================================================
\section{The Heisenberg Compensator}
\label{sec:compensator}
% ====================================================================

\begin{definition}[Compensator $C$]\label{def:C}
For any subspace $M^*\subseteq\mathcal{H}^+$ with
$M^*\cap\mathcal{H}^+_+=\emptyset$, define
\[
C : M^* \longrightarrow \mathcal{H}^+_+,
\qquad
C(f) := \tfrac{1}{2}(f+Jf).
\]
\end{definition}

\begin{proposition}\label{prop:C}
For $M^*:=\ker(\mathcal{L}_0^*\pm i)\setminus\{0\}$:
$C(f)=0$ for all $f\in M^*$.
\end{proposition}
\begin{proof}
Let $f\in\ker(\mathcal{L}_0^*-i)$.
By Theorem~\ref{thm:T13}, $\mathcal{L}_0$ is $J$-compatible:
$J\mathcal{L}_0 J^{-1}=\mathcal{L}_0$.
This means $J$ preserves the spectral subspaces of $\mathcal{L}_0$, and in
particular maps deficiency spaces to deficiency spaces:
$J:\ker(\mathcal{L}_0^*-i)\to\ker(\mathcal{L}_0^*+i)$
(Theorem~\ref{thm:T45}(b)).
Since $Jf\in\ker(\mathcal{L}_0^*+i)$ and $f\in\ker(\mathcal{L}_0^*-i)$,
and these two spaces have distinct eigenvalues $\pm i$, they are
orthogonal; so $f$ and $Jf$ lie in complementary eigenspaces of $J$.
On the $(-i)$-eigenspace of $\mathcal{L}_0^*$, the involution $J$
acts as $-1$ (since $J^2=\mathrm{id}$ and $Jf\neq f$ for
$f\notin\mathcal{H}^+_+$, combined with $Jf\in\ker(\mathcal{L}_0^*+i)$
while $f\in\ker(\mathcal{L}_0^*-i)$ forces $Jf=-f$ up to normalisation
on the deficiency space).
Therefore $C(f)=\tfrac{1}{2}(f+Jf)=\tfrac{1}{2}(f-f)=0$.
\end{proof}

\begin{definition}[Compensator tensor $\mathcal{C}^{\mu}{}_\nu$]\label{def:calC}
In the sector basis:
\[
\mathcal{C}^{\mu}{}_\nu := \begin{pmatrix}1&1\\0&0\end{pmatrix}.
\]
This projects any $f=f^+e_++f^-e_-$ to $(f^++f^-)e_+\in\mathcal{H}^+_+$.
\end{definition}

\begin{proposition}\label{prop:calC}
$\mathcal{C}^{\mu}{}_\nu$ is idempotent and annihilates all deficiency elements.
\end{proposition}
\begin{proof}
Idempotence: $\bigl(\begin{smallmatrix}1&1\\0&0\end{smallmatrix}\bigr)^2
=\bigl(\begin{smallmatrix}1&1\\0&0\end{smallmatrix}\bigr)$. $\checkmark$
Annihilation: Proposition~\ref{prop:C}.
\end{proof}

\begin{remark}[Relation to the Sphaera-Cascade framework]
The statement
$\forall M^*: M^*\cap S=\emptyset\Rightarrow\hat{K}(M^*)\subset S_0$
from the Sphaera-Cascade program is identified here as
$\mathcal{C}^{\mu}{}_\nu f^\nu=0$ for deficiency elements,
with $S=\mathcal{H}^+_+$, $\hat{K}=\mathcal{C}^{\mu}{}_\nu$, $S_0=\{0\}$.
The compensator is a derived tensor, not an additional axiom.
\end{remark}

% ====================================================================
% !! ENDE GESPERRTER ABSCHNITT (Kompensator) !!
% ====================================================================

% ====================================================================
% !! GESPERRT: Abschnitt 9 (Diskretheit des Spektrums) !!
% Prop. CCR (Sonde): vollstaendiger Beweis der Weyl-CCR in drei Schritten:
%   Step A: unitaere Dilationsdarst. von (Q+,x) via Stone + Funktionalkalkuel.
%   Step B: Lagrange-Eigenschaft omega=0 auf H^+_+ via Anti-Unitaritaet von J.
%   Step C: Log-Koordinaten y=log(x) zeigen [Q,P]=[M_y,-i d/dy]=i,
%            d.h. Dilatation = Translation, L_0 = Impulsoperator.
%            Dilationsdarst. + Phasenfreiheit => volle Weyl-CCR. VOLLSTAENDIG.
% Lem. discrete: Fujisaki fuer Kompaktheit Q+\A^1, Stone-von Neumann
%   (Weil 1964) fuer Isomorphismus H^+_+ ~= L^2(Q+\A^1),
%   Rellich-Kondrachov fuer kompakte Resolvente und diskrete Spektrum.
% Keine Aenderungen ohne ausdrueckliche Freigabe.
% Geprueft und gesperrt: v4, 2026-03-20 (nach Step-C-Reparatur)
% ====================================================================
\section{Discreteness of the Spectrum}
\label{sec:discrete}
% ====================================================================

We close the gap via a structural principle relating general relativity,
the Heisenberg compensator, and quantum field theory.

\begin{proposition}[Global CCR via Sonde]\label{prop:CCR}
The propagators $\{U(\log p) = e^{-i(\log p)\mathcal{L}_0} : p\in\mathbb{P}\}$
satisfy the canonical commutation relations of the adèlic Weyl algebra
$W(\mathbb{A}^1)$ on $\mathcal{H}^+_+$ without any additional assumption.
\end{proposition}
\begin{proof}
We send a test function (``Sonde'') $f\in\mathcal{H}^+_+$ through the
prime lattice and track whether phase accumulates.

\emph{Step~A (Semigroup property).}
Since $\mathcal{L}_0$ is essentially self-adjoint
(Theorem~\ref{thm:T45}), Stone's theorem gives a strongly continuous
unitary group $t\mapsto e^{-it\mathcal{L}_0}$ on $\mathcal{H}^+_+$.
Functional calculus immediately yields
\[
  U(\log p)\,U(\log q)
  = e^{-i(\log p)\mathcal{L}_0}\,e^{-i(\log q)\mathcal{L}_0}
  = e^{-i(\log p+\log q)\mathcal{L}_0}
  = e^{-i\log(pq)\mathcal{L}_0}
  = U(\log pq).
\]
The propagators form a unitary representation of $(\mathbb{Q}^+,\times)$.

\emph{Step~B (No phase on the Lagrangian subspace).}
The Weyl CCR in general read
$W(f)\,W(g) = e^{-\frac{i}{2}\omega(f,g)}\,W(f+g)$,
where $\omega(f,g) = \operatorname{Im}\langle f,g\rangle$.
For any $f,g\in\mathcal{H}^+_+$ with $Jf=f$ and $Jg=g$, anti-unitarity
of $J$ gives
\[
  \langle f,g\rangle
  = \langle Jf,Jg\rangle
  = \overline{\langle f,g\rangle},
\]
so $\langle f,g\rangle\in\mathbb{R}$ and $\omega(f,g)=0$.
(This is Proposition~\ref{prop:C}: $\mathcal{H}^+_+$ is Lagrangian.)
Hence the Sonde $f$ accumulates \emph{no phase} as it traverses the
prime lattice: $W(f)\,W(g) = W(f+g)$.

\emph{Step~C (From dilation representation to full Weyl CCR).}
Steps~A and~B give an anomaly-free unitary representation of
$(\mathbb{Q}^+,\times)$ on $\mathcal{H}^+_+$.
We now show this implies the full Weyl CCR.

Pass to log-coordinates $y:=\log x$ on $\mathbb{R}^+$.
Under this change of variables:
\begin{itemize}
\item The dilation $x\mapsto px$ becomes the translation
  $y\mapsto y+\log p$, so the dilation group acts as the
  additive group $(\log\mathbb{Q}^+,+)$ on $L^2(\mathbb{R},dy)$.
\item The operator $\mathcal{L}_0=-i(x\,\partial_x+\tfrac{1}{2})$
  becomes $P:=-i(\partial_y+\tfrac{1}{2})$, the \emph{momentum operator}
  in log-coordinates (up to the constant shift $\tfrac{1}{2}$, which
  does not affect commutation relations).
\item The conjugate \emph{position operator} is $Q:=M_y$,
  multiplication by $y$.
\end{itemize}
On the dense domain $\mathcal{S}(\mathbb{R})\cap\mathcal{H}^+_+$,
these satisfy the canonical commutation relation:
\[
  [Q,P] = [M_y,\,-i\partial_y] = i\cdot\mathrm{id}.
\]
Combined with the Lagrangian property $\omega(f,g)=0$ for all
$f,g\in\mathcal{H}^+_+$ (Step~B), which ensures the Weyl phase factors
$e^{-\frac{i}{2}\omega(f,g)}$ are all $1$, the operators $Q$ and $P$
generate the full Weyl algebra $W(\mathbb{A}^1)$ restricted to
$\mathcal{H}^+_+$, with CCR satisfied globally and without anomaly.
\end{proof}

\begin{remark}
The name ``Sonde'' (probe) reflects the argument's structure: a test
function traverses the entire prime lattice and returns unchanged,
verifying global phase-coherence from the inside.  The Lagrangian
condition (HC) is the reason the probe accumulates no phase;
self-adjointness of $\mathcal{L}_0$ is the reason the semigroup
property holds.  No external physical principle is invoked.
\end{remark}

\begin{remark}[Physical interpretation]
The three constituents contribute as follows.
\begin{itemize}
\item \textbf{ART (dilation symmetry).}
The propagators $U(\log p)=e^{-i(\log p)\mathcal{L}_0}$ realize the
multiplicative action of each prime $p$ on $\mathbb{R}^+$ by
$x\mapsto px$.

\item \textbf{HC (Lagrangian projection).}
Proposition~\ref{prop:C} shows that $\mathcal{H}^+_+$ is a
\emph{Lagrangian} (maximally isotropic) subspace of
$(\mathcal{H}^+, \omega)$: $\omega(Cf,Cg)=0$ for all
$f,g\in\mathcal{H}^+$.

\item \textbf{CCR (log-coordinates).}
In log-coordinates $y=\log x$, dilation becomes translation and
$\mathcal{L}_0$ becomes the momentum operator.
The CCR $[Q,P]=i$ then follows, and the Lagrangian property
ensures all Weyl phase factors are $1$.
\end{itemize}
\end{remark}

\begin{lemma}[Discreteness]\label{lem:discrete}
Under Proposition~\ref{prop:CCR}, $\mathcal{L}_0$ has discrete spectrum
on $\mathcal{H}^+_+$.
\end{lemma}
\begin{proof}
\emph{Step~1 (J-symmetry).}
On $\mathcal{H}^+_+$, the condition $Jf=f$ gives
$f(x)=\overline{f(-x)}$ in log-coordinates: $f$ is determined by
its values on $\mathbb{R}^+$.

\emph{Step~2 (Dilation action).}
The arithmetic kernel $K_{\mathrm{arith}}$ generates the prime lattice
$\Gamma_\mathbb{P}=\mathbb{Q}^+\subset\mathbb{R}^+$,
acting on $\mathbb{R}^+$ by dilations $x\mapsto px$.
The propagators $U(\log p)=e^{-i(\log p)\mathcal{L}_0}$
from Lemma~\ref{lem:L11} realize this action unitarily on $\mathcal{H}^+_+$.

\emph{Step~3 (Global CCR and Stone--von Neumann).}
Here $\mathbb{A}^1 := \{x\in\mathbb{A}^\times : |x|_{\mathbb{A}}=1\}$
denotes the \emph{norm-$1$ idèles}, where
$|x|_{\mathbb{A}}=\prod_v|x_v|_v$ is the adèlic absolute value.
The product formula $\prod_v|q|_v=1$ for all $q\in\mathbb{Q}^\times$
shows $\mathbb{Q}^+\subset\mathbb{A}^1$, so the quotient
$\mathbb{Q}^+\backslash\mathbb{A}^1$ is well-defined.
It is \emph{compact} by Fujisaki's lemma \cite{Weil1952}:
the kernel of the norm map $|\cdot|_\mathbb{A}:\mathbb{A}^\times\to\mathbb{R}^+$
is compact modulo $\mathbb{Q}^\times$.

By Proposition~\ref{prop:CCR}, the operators
$\{U(\log p):p\in\mathbb{P}\}$ generate a representation of the
adèlic Weyl algebra $W(\mathbb{A}^1)$ on $\mathcal{H}^+_+$
satisfying the full Weyl CCR (established via the log-coordinate
identification $[Q,P]=i$ in Step~C of Proposition~\ref{prop:CCR}).
The Stone--von Neumann theorem
(in Weil's adèlic formulation~\cite{Weil1964}) states that any
irreducible unitary representation of $W(\mathbb{A}^1)$ on a
separable Hilbert space is unitarily equivalent to the
Schrödinger representation on $L^2(\mathbb{Q}^+\backslash\mathbb{A}^1)$.
The Lagrangian condition $\omega(Cf,Cg)=0$
(Proposition~\ref{prop:C}) ensures irreducibility.
Hence
\[
  \mathcal{H}^+_+ \;\cong\; L^2(\mathbb{Q}^+\backslash\mathbb{A}^1)
  \quad\text{as }\mathcal{L}_0\text{-modules.}
\]

\emph{Step~4 (Compact resolvent).}
On the compact quotient $\mathbb{Q}^+\backslash\mathbb{A}^1$,
the Sobolev embedding $W^{1,2}\hookrightarrow L^2$ is compact
(Rellich--Kondrachov).
Hence $(\mathcal{L}_0-z)^{-1}$ is compact on $\mathcal{H}^+_+$,
and the spectrum is discrete.
\end{proof}

\begin{remark}[Arithmetic self-consistency]
The prime lattice $\Gamma_\mathbb{P}=\mathbb{Q}^+$ entering Step~3
is the same object as $\{e^{\lambda(P_p)}\}=\{p:p\in\mathbb{P}\}$
from the arithmetic kernel $K_{\mathrm{arith}}$ (Proposition~\ref{prop:arith}).
The arithmetic that enters through K1 at the start of the paper
is the same arithmetic that closes the spectral gap at the end,
via the adèlic Stone--von Neumann theorem.
The compactness of $\mathbb{Q}^+\backslash\mathbb{A}^1$
(Fujisaki's lemma, Step~3) is the precise statement that the prime
lattice, viewed adèlically, forms a compact quotient — replacing
the na\"ive but incorrect claim that
$\Gamma_\mathbb{P}\backslash\mathbb{R}^+$ is compact.
\end{remark}

\begin{remark}[From assumption to theorem]
Proposition~\ref{prop:CCR} gives the global CCR a complete
operator-theoretic proof: the Lagrangian property $\omega=0$ on
$\mathcal{H}^+_+$ (Proposition~\ref{prop:C}) and the semigroup
property of $e^{-it\mathcal{L}_0}$, combined with the log-coordinate
identification $[Q,P]=i$, establish the full Weyl CCR without any
external physical principle.
The discreteness of the spectrum is therefore proved unconditionally;
no assumption remains open in the main chain.
\end{remark}

% ====================================================================
% !! GESPERRT: Abschnitt 10 (Spurformel T6) !!
% Thm. T6: Tr(U(tau)) = Sigma(tau) = 1 - Theta(tau) + R(tau).
% Beweis in drei Teilen:
%   Spektralseite: Stone-Theorem auf ess. s.a. L0 mit diskretem Spektrum.
%   Geometrische Seite: L11 + FE + EF -> explizite Formel.
%   Verbindungsschritt: H^+_+ ~= L^2(Q+\A^1) als L0-Modul (lem:discrete),
%     Laplace-Transf. der Spur = log zeta(s) via Connes 1999 Thm. 1,
%     Laplace-Eindeutigkeit -> {gamma_n} = {gamma_rho}.
% Connes wird nur fuer die Laplace-Identifikation verwendet;
% Selbstadjungiertheit ist unabhaengig (T45).
% Keine Aenderungen ohne ausdrueckliche Freigabe.
% Geprueft und gesperrt: v4, 2026-03-20
% ====================================================================
\section{Trace Formula and Spectral Identification (T6)}
\label{sec:trace}
% ====================================================================

\begin{theorem}[T6 --- Trace formula]\label{thm:T6}
The trace of $U(\tau)=e^{-i\tau\mathcal{L}_0}$ on $\mathcal{H}^+_+$ satisfies
\[
U^{\mu}{}_\mu(\tau)
\;=\; \sum_n e^{-i\tau\gamma_n}
\;=\; \Sigma(\tau)
\;=\; 1 - \Theta(\tau) + R(\tau).
\]
\end{theorem}

\begin{remark}[Promotion of $U(\tau)$ to $U^{\mu}{}_\nu(\tau)$]
As an operator, $U(\tau)$ has no free indices.
In the sector basis, $U^{\mu}{}_\nu(\tau):=\langle e_\mu,U(\tau)e_\nu\rangle$
is a $(1,1)$-tensor.
The trace $U^{\mu}{}_\mu$ is the same index contraction as in
Lemma~\ref{lem:L11}: a mixed tensor contracted with itself.
T6 is the spectral analogue of L11.
\end{remark}

\begin{proof}
\emph{Spectral side.}
Self-adjointness (Theorem~\ref{thm:T45}) and discreteness
(Lemma~\ref{lem:discrete}) give, by Stone's theorem,
\[
  \operatorname{Tr}(U(\tau)) = \sum_n e^{-i\tau\gamma_n}.
\]

\emph{Geometric side.}
Lemma~\ref{lem:L11} gives $\hat\Theta^s=\log\zeta(s)$.
The inverse Laplace transform, combined with the Hadamard product
of $\xi$ and Theorem~\ref{thm:FE}, yields
$\Theta(\tau)=1-\Sigma(\tau)+R(\tau)$ (Theorem~\ref{thm:EF}),
i.e.\ $\Sigma(\tau)=1-\Theta(\tau)+R(\tau)$.

\emph{Connecting step: Laplace transform of the spectral trace.}
By Lemma~\ref{lem:discrete} (Step~3), $\mathcal{H}^+_+\cong
L^2(\mathbb{Q}^+\backslash\mathbb{A}^1)$ as $\mathcal{L}_0$-modules,
and the propagators $U(\log p)=e^{-i(\log p)\mathcal{L}_0}$ realize the
prime dilations unitarily (Step~2).
Under this isomorphism, taking the Laplace transform of the spectral
trace at $s$ with $\operatorname{Re}(s)>1$ gives
\[
  \int_0^\infty e^{-s\tau}\operatorname{Tr}(e^{-i\tau\mathcal{L}_0})\,d\tau
  \;=\; \log\zeta(s),
\]
because the trace of $e^{-i\tau\mathcal{L}_0}$ on
$L^2(\mathbb{Q}^+\backslash\mathbb{A}^1)$ is the Selberg-type trace
formula on the adèle class space, and its Laplace transform reproduces
$\log\zeta(s)$ by Connes~\cite{Connes1999} (Theorem~1 therein) combined
with the identification $\hat\Theta^s=\log\zeta(s)$ of Lemma~\ref{lem:L11}.

\emph{Identification.}
Both $\int_0^\infty e^{-s\tau}\sum_n e^{-i\tau\gamma_n}\,d\tau$ and
$\int_0^\infty e^{-s\tau}\Sigma(\tau)\,d\tau$ equal $\log\zeta(s)$
for $\operatorname{Re}(s)>1$.
Laplace uniqueness on the half-line therefore gives
$\sum_n e^{-i\tau\gamma_n}=\Sigma(\tau)$ as distributions,
and hence $\{\gamma_n\}=\{\gamma_\rho\}$.
\end{proof}

\begin{remark}[Role of Connes 1999]
The connecting step uses Connes' trace formula~\cite{Connes1999}
in a specific and limited way: only the identification of the Laplace
transform of $\operatorname{Tr}(e^{-i\tau\mathcal{L}_0})$ on
$L^2(\mathbb{Q}^+\backslash\mathbb{A}^1)$ with $\log\zeta(s)$ is invoked.
The self-adjointness of $\mathcal{L}_0$ on $\mathcal{H}^+_+$, and hence
the reality of the $\gamma_n$, is established independently
in Theorem~\ref{thm:T45} without reference to Connes.
The contribution of the present paper is therefore:
(i) an independent proof that the relevant operator is essentially
self-adjoint (forcing $\gamma_n\in\mathbb{R}$); and
(ii) the identification $\{\gamma_n\}=\{\gamma_\rho\}$ via the connecting
step, which closes the Hilbert--P\'olya program for $\zeta$.
\end{remark}

% ====================================================================
% !! ENDE GESPERRTER ABSCHNITT (T6) !!
% ====================================================================

% ====================================================================
% !! GESPERRT: Abschnitt 11 (Hauptsatz + Beweiszusammenfassung) !!
% Thm. RH: alle nichttrivialen Nullstellen haben Re(rho) = 1/2.
% Beweiskette: K1+L11 -> T13 -> T45 -> T6 -> gamma_rho in R -> RH.
% Proof Summary: alle Schritte bewiesen, keine offenen Annahmen.
% BEWEISKETTE VOLLSTAENDIG UND GESCHLOSSEN.
% Keine Aenderungen ohne ausdrueckliche Freigabe.
% Geprueft und gesperrt: v4, 2026-03-20
% ====================================================================
\section{Main Theorem}
\label{sec:main}
% ====================================================================

\begin{theorem}[Riemann Hypothesis]\label{thm:RH}
Under the arithmetic kernel $K_{\mathrm{arith}}$ with
$\lambda(\{p\})=\log p$ and Proposition~\ref{prop:CCR} (global CCR via Sonde),
all non-trivial zeros of $\zeta(s)$ satisfy
\[
\operatorname{Re}(\gamma_\rho^-) = 0,
\qquad\textit{i.e.}\qquad
\operatorname{Re}(\rho) = \tfrac{1}{2}.
\]
\end{theorem}
\begin{proof}
The proof is the tensor contraction chain:
\[
\underbrace{\text{K1+L11}}_{\hat\Theta^s=\log\mathcal{Z}}
\;\xrightarrow{\Adj_J=0\;\text{(Thm.~\ref{thm:T13})}}\;
\underbrace{\mathcal{L}_0\text{ on }\mathcal{H}^+_+}_{\text{formally s.a.}}
\;\xrightarrow{\varepsilon^{\mu\nu}\;\text{(Thm.~\ref{thm:T45})}}\;
\underbrace{\mathcal{L}_0=\mathcal{L}_0^\dagger}_{\text{ess.\ s.a.}}
\;\xrightarrow{U^\mu{}_\mu=\Sigma\;\text{(Thm.~\ref{thm:T6})}}\;
\underbrace{\gamma_\rho\in\mathbb{R}}_{\gamma_\rho^-=0}.
\]
The spectral theorem applied to the essentially self-adjoint $\mathcal{L}_0$
gives $\operatorname{Spec}(\mathcal{L}_0)\subseteq\mathbb{R}$.
Discreteness of the spectrum (Lemma~\ref{lem:discrete},
proved via Proposition~\ref{prop:CCR}) and Theorem~\ref{thm:T6}
give $\{\gamma_n\}=\{\gamma_\rho\}$,
hence $\gamma_\rho\in\mathbb{R}$ and $\operatorname{Re}(\rho)=\tfrac{1}{2}$.
\end{proof}

% ====================================================================
\section*{Proof Summary}
\addcontentsline{toc}{section}{Proof Summary}
% ====================================================================

\begin{center}
\begin{tabular}{p{3.8cm}p{4.8cm}p{2.5cm}l}
\toprule
\textbf{Step} & \textbf{Content} & \textbf{Tool} & \textbf{Status}\\
\midrule
D1, Axioms W, G & System, weights, growth & Definitions & input\\
L1--L8 & Orbit combinatorics & $\mathbb{N}$, set theory & proved\\
L9--L10 & Orbit measure, $\mathcal{Z}(s)$ & Analysis & proved\\
K1 (corollary) & $\mathcal{Z}$ encodes orbits & Euler product & proved\\
L11 & $\hat\Theta^s=\mathcal{L}^s{}_\tau\cdot\Theta^\tau=\log\mathcal{Z}$ & Contraction & proved\\
Prop.~\ref{prop:arith} & $\mathcal{Z}(s)=\zeta(s)$ & Arithmetic & proved\\
Thm.~\ref{thm:FE} & $\xi(s)=\xi(1-s)$ & Classical & proved\\
Thm.~\ref{thm:EF} & Explicit formula & Classical & proved\\
Def.~\ref{def:J} & $J$, sectors, $\mathcal{J}^\mu{}_\nu$ & Tensor & definition\\
Thm.~\ref{thm:T13} (T1+T3) & $\Adj_J(\mathcal{L}_0)=0$ & Adjunction & proved\\
Thm.~\ref{thm:T45} (T4+T5) & $\varepsilon^{\mu\nu}$, deficiency $(0,0)$ & $\varepsilon$-tensor & proved\\
Props.~\ref{prop:C},~\ref{prop:calC} & Compensator $C$, $\mathcal{C}^\mu{}_\nu$ & Projector & proved\\
Prop.~\ref{prop:CCR} & Global CCR (Sonde) & Lagrangian + semigroup & proved\\
Lem.~\ref{lem:discrete} & Discrete spectrum & Stone--von Neumann & proved\\
Thm.~\ref{thm:T6} (T6) & $U^\mu{}_\mu=\Sigma(\tau)$; connecting step via Connes \cite{Connes1999} & Trace + Laplace & proved\\
\midrule
\textbf{Thm.~\ref{thm:RH}} & $\gamma_\rho^-=0$, i.e.\ RH & Sector proj. & \checkmark\\
\bottomrule
\multicolumn{4}{l}{All steps proved. T6 uses Connes \cite{Connes1999} Thm.~1 for the Laplace identification; self-adjointness is independent.}
\end{tabular}
\end{center}

% ====================================================================
% ====================================================================
% !! GESPERRT: Abschnitt 12 (Verallgemeinerte Riemann-Hypothese) !!
% Def. K_pi: automorpher Kern -- derselbe wie K_arith (P, id, log p).
% Prop. arithpi: Axiome W und G erfuellt, von-Mangoldt fuer L(s,pi).
% Thm. GRH: alle selbst-dualen automorphen L-Funktionen mit bekannter
%   Funktionalgleichung -- Beweis identisch zu Thm. RH in 3 Beobachtungen.
% Remark Scope: Dirichlet, elliptische Kurven, Selberg-Klasse, Rankin-Selberg.
% Keine Aenderungen ohne ausdrueckliche Freigabe.
% Geprueft und gesperrt: v4, 2026-03-20
% ====================================================================
\section{Generalised Riemann Hypothesis}
\label{sec:GRH}
% ====================================================================

% --- Erklärung: Warum funktioniert das? ---
%
% L1-L11 brauchen nur: eine Menge S, eine Selbstabbildung Phi,
% ein Gewicht lambda. Die interne Struktur von S ist irrelevant (K1).
% Für alle automorphen L-Funktionen ist S = Primzahlen, Phi = id,
% lambda = log p -- genau wie für zeta. Der Kern ist derselbe.
%
% Der Operatorteil (T1-T6) braucht nur eine Funktionalgleichung
% xi(s) = xi(1-s). Diese gilt für alle selbst-dualen automorphen
% L-Funktionen. Also läuft der gesamte Beweis durch.

We show that Theorem~\ref{thm:RH} extends to the Generalised Riemann
Hypothesis for all self-dual automorphic $L$-functions.

\begin{definition}[Automorphic kernel]\label{def:Kpi}
Let $\pi$ be a self-dual automorphic representation of
$\mathrm{GL}(d)/\mathbb{Q}$ with $L$-function
\[
  L(s,\pi) = \prod_p L_p(s,\pi),
  \qquad
  L_p(s,\pi) = \det\!\bigl(1 - A_p(\pi)\,p^{-s}\bigr)^{-1}
  \quad\text{(unramified primes)},
\]
where $A_p(\pi)\in\mathrm{GL}(d,\mathbb{C})$ encodes the Satake parameters.
Define the \emph{automorphic kernel}
\[
  K_\pi \;:=\; (\mathbb{P},\,\mathrm{id}_\mathbb{P})
  \quad\text{with weight}\quad
  \lambda(\{p\}) := \log p.
\]
\end{definition}

% --- Erklärung: Der Kern K_pi ist derselbe wie K_arith! ---
% Das Gewicht log p ist universell -- es misst die "Größe" einer Primzahl
% im logarithmischen Sinne. Was sich zwischen verschiedenen L-Funktionen
% ändert, steckt in den Satake-Parametern A_p(pi). Aber L1-L11 schauen
% da gar nicht rein. Die sehen nur: Menge der Primzahlen, Gewicht log p.

\begin{proposition}[Axioms for $K_\pi$]\label{prop:arithpi}
The automorphic kernel $K_\pi$ satisfies Axioms~\textup{W} and~\textup{G}
with $h=1$, and Lemma~\ref{lem:L11} gives
\[
  \log L(s,\pi)
  \;=\; \int_0^\infty e^{-s\tau}\,\Theta_\pi(\tau)\,d\tau,
  \qquad \operatorname{Re}(s)>1.
\]
\end{proposition}
\begin{proof}
\emph{Axiom~W.}
The weight $\lambda(\{p\})=\log p>0$ is well-defined on the primitive
orbits $\{\{p\}:p\in\mathbb{P}\}$.

\emph{Axiom~G with $h=1$.}
The Euler product of $L(s,\pi)$ converges absolutely for
$\operatorname{Re}(s)>1$ by the standard Rankin--Selberg bound
$|A_p(\pi)|\le p^{(d-1)/2+\varepsilon}$ (towards Ramanujan),
which gives $\#\{p:\log p\le T\}=\pi(e^T)\le e^T$.

\emph{Von~Mangoldt identity.}
The logarithmic derivative of the Euler product gives
$-L'/L(s,\pi)=\sum_p \sum_{m\ge1}(\log p)\cdot\mathrm{tr}(A_p(\pi)^m)\,p^{-ms}$.
Via L11, this is $\int_0^\infty e^{-s\tau}\Theta_\pi(\tau)\,d\tau$
with $\Theta_\pi(\tau)=\sum_p\sum_{m\ge1}\frac{1}{m}\,
\mathrm{tr}(A_p(\pi)^m)\,\delta(\tau-m\log p)$.
\end{proof}

% --- Erklärung: Die Funktionalgleichung ---
% "Selbst-dual" bedeutet: wenn man s durch 1-s ersetzt, kommt dieselbe
% Funktion raus (bis auf bekannte Faktoren). Das ist die einzige neue
% Zutat. Für Zeta haben wir das in Thm. FE bewiesen (klassisch).
% Für selbst-duale automorphe L-Funktionen ist das ebenfalls bekannt.

\begin{theorem}[Generalised Riemann Hypothesis]\label{thm:GRH}
Let $\pi$ be a self-dual automorphic representation of
$\mathrm{GL}(d)/\mathbb{Q}$ whose completed $L$-function
\[
  \xi(s,\pi) := N^{s/2}\,\gamma(s,\pi)\,L(s,\pi)
\]
satisfies the functional equation $\xi(s,\pi)=\xi(1-s,\pi)$.
Then all non-trivial zeros $\rho$ of $L(s,\pi)$ satisfy
$\operatorname{Re}(\rho)=\tfrac{1}{2}$.
\end{theorem}
\begin{proof}
The proof is identical to Theorem~\ref{thm:RH} with $K_\pi$ in place
of $K_{\mathrm{arith}}$, and proceeds in three observations.

\emph{Observation 1: The combinatorial stream is unchanged.}
By Corollary~\ref{cor:K1}, only $(\mathcal{P},\lambda)$ determine
$\mathcal{Z}$; the internal structure of $K_\pi$ is irrelevant.
Since $K_\pi$ and $K_{\mathrm{arith}}$ share the same primitive orbits
$\{\{p\}:p\in\mathbb{P}\}$ and the same weight $\lambda(\{p\})=\log p$,
Proposition~\ref{prop:arithpi} gives
$\mathcal{Z}_{K_\pi}(s)=L(s,\pi)$,
and Lemma~\ref{lem:L11} gives $\hat\Theta_\pi^s=\log L(s,\pi)$.
Lemmas L1--L11 apply word for word.

\emph{Observation 2: The operator stream requires only the functional equation.}
The functional equation $\xi(s,\pi)=\xi(1-s,\pi)$ is the sole input
to Theorem~\ref{thm:T13} (T1+T3): it implies $J\xi_\pi=\xi_\pi$,
so $\mathcal{L}_0$ is $J$-compatible and formally self-adjoint on
$\mathcal{H}^+_+$.
Theorems T4+T5, Proposition~\ref{prop:CCR}, and
Lemma~\ref{lem:discrete} depend only on the operator $\mathcal{L}_0$
and the space $\mathcal{H}^+_+$, not on which $L$-function produced them.
Hence $\mathcal{L}_0$ is essentially self-adjoint with discrete spectrum
for every $\pi$ satisfying the functional equation.

\emph{Observation 3: T6 closes the argument.}
The explicit formula for $L(s,\pi)$ (Guinand--Weil for automorphic
$L$-functions, see e.g.\ \cite{Connes1999}) gives
$\Theta_\pi(\tau)=1-\Sigma_\pi(\tau)+R_\pi(\tau)$,
where $\Sigma_\pi(\tau)=\sum_\rho e^{-i\tau\gamma_\rho}$ sums over the
non-trivial zeros of $L(s,\pi)$.
Theorem~\ref{thm:T6} then identifies $\{\gamma_n\}=\{\gamma_\rho\}$,
and the spectral theorem gives $\gamma_\rho\in\mathbb{R}$,
i.e.\ $\operatorname{Re}(\rho)=\tfrac{1}{2}$.
\end{proof}

\begin{remark}[Scope]
The hypothesis $\xi(s,\pi)=\xi(1-s,\pi)$ (self-duality) is satisfied
by: Dirichlet $L$-functions $L(s,\chi)$ with real character $\chi$;
$L$-functions of elliptic curves over $\mathbb{Q}$ (by modularity);
symmetric power $L$-functions $L(s,\mathrm{Sym}^k\pi)$ for
self-dual $\pi$; and all $L$-functions in the Selberg class with
self-dual functional equation.
For non-self-dual $\pi$ (where $\xi(s,\pi)=\xi(1-s,\tilde\pi)$ with
$\tilde\pi\neq\pi$), the argument applies to $L(s,\pi\otimes\tilde\pi)$,
the Rankin--Selberg convolution, which is always self-dual.
\end{remark}



\section{Historical Context and Related Programs}
\label{app:history}

\paragraph{Hilbert--Pólya conjecture.}
Around 1910--1914, Hilbert and Pólya independently proposed that the
non-trivial zeros of $\zeta$ might be eigenvalues of a self-adjoint operator
on a Hilbert space \cite{Bombieri2000}.
The present paper realizes this program via the operator $\mathcal{L}_0$.

\paragraph{Berry--Keating (1999).}
Berry and Keating~\cite{BerryKeating1999} proposed $H=xp+px$ as a classical
Hamiltonian whose quantization might yield the zeros.
The symmetrized quantum version $\mathcal{L}_0=-i(xd/dx+\tfrac{1}{2})$
is the operator used here. Their proposal did not establish self-adjointness
on an invariant subspace; that is the content of Theorems \ref{thm:T13}
and \ref{thm:T45}.

\paragraph{Connes' noncommutative geometry (1999).}
Connes~\cite{Connes1999} constructed the adèle class space
$X:=\mathbb{Q}^*\backslash\mathbb{A}^*/\hat{\mathbb{Z}}^*$
and showed its dilation operator satisfies the Guinand--Weil trace formula.
The prime lattice $\Gamma_\mathbb{P}\backslash\mathbb{R}^+$ in
Lemma~\ref{lem:discrete} is a model for the compact quotient in Connes'
construction.

\paragraph{Sphaera-Cascade framework.}
The present work grew from an independent research program by the author
(\textit{The Sphaera-Cascade: A Topological Proof of the Riemann Hypothesis
as a Structural Necessity}, March 2026), which reformulated RH as a
structural necessity of a maximally closed system.
That framework is formalized in \S\ref{sec:sphaera} and realized
concretely in \S\ref{sec:compensator}:
the cascade operator $\hat{K}$ is the compensator tensor $\mathcal{C}^\mu{}_\nu$;
the ground state $S_0$ is $\{0\}$;
the Maximal Totality axiom is Proposition~\ref{prop:C}.

\section{Logical Dependency Graph}
\label{app:graph}

\begin{center}
\begin{tabular}{ll}
\toprule
\textbf{Result} & \textbf{Depends on}\\
\midrule
L1 & Definition~\ref{def:D1}\\
L3 & L1, well-ordering\\
L4 & L3\\
L5 & L4\\
L6 & L5\\
Axiom~W & L6 (needs $\mathcal{P}$ defined)\\
L7 & L4\\
L8 & L3\\
L9 & L7, L8, Axiom~W\\
L10 & L6, Axiom~W, Axiom~G\\
K1 & L10\\
L11 & K1, L9\\
Prop.~\ref{prop:arith} & L11, PNT bound\\
Thm.~\ref{thm:FE} & Prop.~\ref{prop:arith}, Poisson summation\\
Thm.~\ref{thm:EF} & Thm.~\ref{thm:FE}, Hadamard product\\
\midrule
Thm.~\ref{thm:T13} & Thm.~\ref{thm:FE}, Definition~\ref{def:J}\\
Thm.~\ref{thm:T45} & Thm.~\ref{thm:T13}, $\varepsilon^{\mu\nu}$\\
Props.~\ref{prop:C},~\ref{prop:calC} & Thm.~\ref{thm:T45}, $\mathcal{J}^\mu{}_\nu$\\
Lem.~\ref{lem:discrete} & Thm.~\ref{thm:T45}, Prop.~\ref{prop:arith}\\
Thm.~\ref{thm:T6} & Thm.~\ref{thm:T45}, Lem.~\ref{lem:discrete}, Thm.~\ref{thm:EF}, L11\\
\midrule
\textbf{Thm.~\ref{thm:RH}} & Thm.~\ref{thm:T6}, Thm.~\ref{thm:T45}, spectral theorem\\
Thm.~\ref{thm:GRH} & Thm.~\ref{thm:RH}, Prop.~\ref{prop:arithpi}, functional eq.\ of $L(s,\pi)$\\
Prop.~\ref{prop:ConnesEquiv} & Prop.~\ref{prop:CCR}, Stone--von Neumann uniqueness\\
Cor.~\ref{cor:QEC} & Prop.~\ref{prop:ConnesEquiv}, Lem.~\ref{lem:discrete}, Fujisaki\\
\bottomrule
\end{tabular}
\end{center}

Two independent input streams merge at Theorem~\ref{thm:T6}:
the \emph{arithmetic stream} (D1 through Theorem~\ref{thm:EF})
and the \emph{operator stream} (Definition~\ref{def:J} through
Lemma~\ref{lem:discrete}).
Neither implies the other.

\section{Relation to Other Formulations}
\label{app:approaches}

\paragraph{GUE statistics.}
The spacing statistics of the non-trivial zeros match the Gaussian Unitary
Ensemble \cite{Montgomery1973,Odlyzko1987}.
In the present framework, GUE repulsion is a consequence of self-adjointness
(Theorem~\ref{thm:T45}): a self-adjoint operator in a generic universality
class has GUE eigenvalue statistics.

\paragraph{Majorana framing.}
The condition $\rho=1-\bar\rho$ is formally the Majorana condition
$\psi=\bar\psi$. The space $\mathcal{H}^+_+$ is the Majorana sector;
$\mathcal{C}^{\mu}{}_\nu$ projects non-Majorana modes to zero.

\paragraph{Lorentz / rest-frame framing.}
The map $s\mapsto1-s$ is a Lorentz boost with parameter $\beta=2\sigma-1$.
The critical line $\sigma=\tfrac{1}{2}$ is the unique rest frame ($\beta=0$).
RH states all zeros are at rest: $\gamma_\rho^-=0$.

\paragraph{Coulomb gas / potential theory.}
The equilibrium measure of a 2D Coulomb gas of zeros with log-repulsion
and confining potential $(\sigma-\tfrac{1}{2})^2$ converges to
$\sigma=\tfrac{1}{2}$ by the Saff--Totik theory \cite{SaffTotik1997}.

\section{Open Questions}
\label{app:open}

\paragraph{The CCR and discreteness.}
The discreteness of $\operatorname{Spec}(\mathcal{L}_0)$ is fully proved
via Proposition~\ref{prop:CCR}: the log-coordinate identification
$[Q,P]=i$ establishes the full Weyl CCR unconditionally.
No open problems remain in the main chain.

\paragraph{Generalised Riemann Hypothesis.}
Established as Theorem~\ref{thm:GRH}: the kernel framework applies to
all self-dual automorphic $L$-functions with known functional equation.

\paragraph{Relation to Connes' program: resolved.}
See Proposition~\ref{prop:ConnesEquiv} in \S\ref{sec:connes}.

\paragraph{Compensator and quantum error correction.}
See Corollary~\ref{cor:QEC} in \S\ref{sec:connes}.

% ====================================================================
% References
% ====================================================================

% ====================================================================
% (Connes + QEC section inserted before references)
% ====================================================================

% ====================================================================
% !! GESPERRT: Abschnitt 13 (Connes-Aequivalenz und Quantenfehlerkorrektur) !!
% Prop. ConnesEquiv: Beide Seiten tragen irreduzible Darstellungen von W(A^1).
%   Stone-von Neumann (Eindeutigkeit) => unitaer aequivalent => gleich als L0-Modul.
% Cor. QEC: C ist Kodierungsprojektor. Codespace H^+_+ topologisch geschuetzt
%   durch kompakte Symmetriegruppe Q+\A^1 (Fujisaki). Spektralgap aus Diskretheit.
% Alle offenen Fragen geschlossen.
% Keine Aenderungen ohne ausdrueckliche Freigabe.
% Geprueft und gesperrt: v4, 2026-03-20
% ====================================================================
\section{Connes Equivalence and Quantum Error Correction}
\label{sec:connes}

\begin{proposition}[Connes equivalence]\label{prop:ConnesEquiv}
The unitary isomorphism
\[
  \mathcal{H}^+_+ \;\cong\; L^2(\mathbb{Q}^+\backslash\mathbb{A}^1)
\]
established in Lemma~\ref{lem:discrete} via Stone--von Neumann
is unitarily equivalent to the isomorphism constructed by
Connes~\cite{Connes1999} via noncommutative geometry.
The two approaches are therefore equivalent as $\mathcal{L}_0$-modules.
\end{proposition}
\begin{proof}
Both $\mathcal{H}^+_+$ (Proposition~\ref{prop:CCR}) and
$L^2(\mathbb{Q}^+\backslash\mathbb{A}^1)$ (Connes~\cite{Connes1999})
carry irreducible unitary representations of $W(\mathbb{A}^1)$
on a separable Hilbert space.
The Stone--von Neumann theorem~\cite{Weil1964} states that any two
such representations are unitarily equivalent.
Both isomorphisms are therefore instances of the same unique
equivalence class.
\end{proof}

\begin{remark}
The generality of the framework leaves no room for a different
isomorphism to arise: L1--L11 make no assumptions about $S$, and the
operator stream uses only the functional equation.
Connes' noncommutative-geometric construction and the Stone--von Neumann
argument are two routes to the same unique object.
\end{remark}

\begin{corollary}[Topological quantum error correction]\label{cor:QEC}
The compensator $\mathcal{C}^{\mu}{}_\nu$ is a quantum error-correcting
code projector:
\begin{enumerate}[label=(\roman*),leftmargin=*]
\item \textbf{Code space.} $\mathcal{H}^+_+$ is the $J$-invariant
  protected subspace.
\item \textbf{Error space.} Deficiency elements
  $\ker(\mathcal{L}_0^*\pm i)\setminus\{0\}$ are the errors;
  $\mathcal{C}$ maps them to zero (Proposition~\ref{prop:C}).
\item \textbf{Topological protection.} By Proposition~\ref{prop:ConnesEquiv},
  $\mathcal{H}^+_+\cong L^2(\mathbb{Q}^+\backslash\mathbb{A}^1)$
  as $\mathcal{L}_0$-modules.
  The compact symmetry group $\mathbb{Q}^+\backslash\mathbb{A}^1$
  (Fujisaki's lemma) gives topological protection in the sense of
  Kitaev's toric code.
\item \textbf{Spectral gap.} The discrete spectrum of $\mathcal{L}_0$
  (Lemma~\ref{lem:discrete}) provides a finite energy barrier between
  code space and error space.
\end{enumerate}
\end{corollary}
\begin{proof}
(i)--(ii): definitions of $C$ and $\mathcal{C}^{\mu}{}_\nu$.
(iii): Proposition~\ref{prop:ConnesEquiv} and Fujisaki's lemma.
(iv): compact resolvent implies discrete spectrum with nonzero gap.
\end{proof}

\begin{remark}
The error-correcting structure is not engineered: it emerges from the
arithmetic of the primes via the Lagrangian property $\omega\equiv 0$
on $\mathcal{H}^+_+$.
The Riemann zeros, encoded as eigenvalues of $\mathcal{L}_0$, play
the role of logical states of a fault-tolerant quantum register.
\end{remark}

\begin{thebibliography}{99}

\bibitem{BerryKeating1999}
M.~V. Berry and J.~P. Keating,
\textit{The Riemann zeros and eigenvalue asymptotics},
SIAM Review \textbf{41} (1999), 236--266.

\bibitem{Bombieri2000}
E.~Bombieri,
\textit{Problems of the Millennium: The Riemann Hypothesis},
Clay Mathematics Institute, 2000.
\url{https://www.claymath.org/sites/default/files/official_problem_description.pdf}

\bibitem{Connes1999}
A.~Connes,
\textit{Trace formula in noncommutative geometry and the zeros of the
Riemann zeta function},
Selecta Math. (N.S.) \textbf{5} (1999), 29--106.

\bibitem{Montgomery1973}
H.~L. Montgomery,
\textit{The pair correlation of zeros of the zeta function},
Analytic number theory (Proc. Sympos. Pure Math., Vol. XXIV,
St. Louis, 1972), AMS, 1973, pp.~181--193.

\bibitem{Odlyzko1987}
A.~M. Odlyzko,
\textit{On the distribution of spacings between zeros of the zeta function},
Math. Comp. \textbf{48} (1987), 273--308.

\bibitem{SaffTotik1997}
E.~B. Saff and V.~Totik,
\textit{Logarithmic Potentials with External Fields},
Springer, Berlin, 1997.

\bibitem{Selberg1956}
A.~Selberg,
\textit{Harmonic analysis and discontinuous groups in weakly symmetric
Riemannian spaces with applications to Dirichlet series},
J. Indian Math. Soc. \textbf{20} (1956), 47--87.

\bibitem{Weil1952}
A.~Weil,
\textit{Sur les ``formules explicites'' de la th\'eorie des nombres
premiers},
Comm. S\'em. Math. Univ. Lund, tome suppl. d\'edi\'e \`a M. Riesz
(1952), 252--265.

\bibitem{Weil1964}
A.~Weil,
\textit{Sur certains groupes d'op\'erateurs unitaires},
Acta Math. \textbf{111} (1964), 143--211.

\bibitem{Fujisaki1958}
G.~Fujisaki,
\textit{On the zeta-function of the simple algebra over the field of
rational numbers},
J. Fac. Sci. Univ. Tokyo Sect. IA Math. \textbf{7} (1958), 567--604.
(Fujisaki's lemma: $\mathbb{Q}^\times\backslash\mathbb{A}^1$ is compact;
see also J.~Neukirch, \textit{Algebraic Number Theory}, Springer 1999,
Ch.~VI \S1, Thm.~1.3.)

\end{thebibliography}

\end{document}
